% ============================================================
%  Harald Høffding, "Filosofiske Problemer" (Copenhagen, 1902)
%  ENGLISH TRANSLATION
%
%  Translation-only project, on the model of danske-filosoffer and
%  kierkegaard-filosof: a clean modern e-book edition exists (SAGA Egmont /
%  Lindhardt & Ringhof, ISBN 9788726362992), so no Danish transcription is
%  produced and none is published. The Danish text is the publisher's
%  copyrighted edition; only this English translation is ours to publish.
%
%  Source of truth for translating:
%    ~/bibliotek/Høffding, Harald/Filosofiske_problemer.epub
%    (© 1902, 2020 Lindhardt og Ringhof; ISBN 9788726362992)
%  Chapter text in OPS/s003–s007; the 63 endnotes in OPS/s010-Notes-01.xhtml,
%  keyed rw-num-note-N and restored here as \footnote{} at their anchors.
%
%  PAGE MAP — READ OFF THE 1902 PRINT.  The Royal Danish Library scan of the
%  first edition is at https://www.kb.dk/e-mat/dod/11140400794C-bw.pdf
%  (local copy: ~/bibliotek/Høffding, Harald/1902-filosofiske-problemer.pdf;
%  90 printed pages; PDF page = printed page + 13).  The map below is taken
%  from the scan's own running heads and section openings, not inferred:
%    Indledning  pp.  1–5
%    I.   pp.  6–27   §1  6   §2 12   §3 22   §4 25
%    II.  pp. 28–52   §1 28   §2 31   §3 39   §4 43   §5 48
%    III. pp. 53–69   §1 53   §2 57   §3 59   §4 62   §5 65
%    IV.  pp. 70–84   §1 70   §2 72   §3 75   §4 77   §5 78   §6 80   §7 82
%         A. Det etiske Problem p. 72;  B. Det religiøse Problem p. 78.
%    Notes printed pp. 85–90.
%  This SUPERSEDES two earlier maps, both built from the endnotes' own "(S. N)"
%  refs before the scan was available — one that read II as 28–53 / III 54–70,
%  and a least-squares correction that read II as 28–51 / III 52–69.  The
%  regression was close (it got IV's opening exactly, and III to within one
%  page) but it is no longer needed and should not be relied on.
%
%  THE E-BOOK RENUMBERED THE NOTES.  The 1902 print numbers its endnotes
%  1–60; the SAGA e-book numbers them 1–63.  The offset opens gradually (0 at
%  note 4, −1 by note 29, −3 from note 41 on), so e-book note N is print note
%  N−3 in the later chapters.  Every note's own "(S. N)" page ref is identical
%  in both, so the renumbering did not disturb the page map — but it DID break
%  the book's two internal note cross-references, which the e-book left
%  pointing at their old numbers:
%    e-book note 36 "Se Note 30" — print note 33 → print note 30 = Maxwell /
%      Hertz.  In e-book numbering that note is 31, not 30 (30 is Mach).
%    e-book note 40 "Se ovenfor Note 24" — print note 37 → print note 24 =
%      Ostwald's Naturphilosophie.  In e-book numbering that note is 25.
%  This translation numbers its footnotes 1–63 in e-book order, so both
%  references are emended to that sequence: "See Note 31" and "See Note 25".
%  An earlier note here called "Se Note 30" a probable authorial slip.  It is
%  not: it is correct in the print, and only the e-book's renumbering broke it.
%
%  GREEK.  One Greek quotation only, in the last footnote of the book (note 63,
%  Plato, Phaedrus 249 C).  The 1902 print sets it in period Greek types, with
%  the variant rho and theta forms; the e-book preserved those variants but
%  corrupted the text, substituting Latin o and accented Latin vowels.  Read
%  off the page image and set here in standard Greek letterforms with the
%  correct accents: πρὸς οἷσπερ θεὸς ὢν θεῖός ἐστι.  Normalising the variant
%  letterforms is a typographic decision, not a textual one, on a par with not
%  reproducing the long s.
%  This file therefore needs \usepackage{textalpha} to build — as it already
%  did for the α/β of II §2.  THE SANDBOX RECIPE STRIPS textalpha AND REPLACES
%  GREEK WITH "[Gr]", so a clean sandbox compile does NOT prove the Greek
%  renders.  Check note 63 and p. 19 on the first local build.
%
%  NOTE 43's "(S. 41)" IS A MISPRINT, as suspected, and it is in the 1902
%  print too (print note 40) — not an e-book defect.  Its neighbours read
%  (S. 44) and (S. 45), and II §4 runs pp. 43–47, so the Bradley–Bosanquet
%  passage falls on p. 44 or 45.  Left as printed; recorded here.
%
%  TITLE.  Rendered "Philosophical Problems", the literal title and the one
%  Fisher's 1905 Macmillan translation uses in its running heads, rather than
%  its title-page "The Problems of Philosophy" — which collides with Russell
%  (1912) and misstates a book that is about four problems, not the subject at
%  large.
%
%  INTERNAL CROSS-REFERENCES.  Høffding cites his own sections as "(II, 3)",
%  "(se III, 2)" and so on.  EIGHT such references occur, and SEVEN of them
%  name a chapter one higher than the printed chapter numeral — as though the
%  unnumbered Indledning were counted as chapter I.  Each was checked against
%  the section it must mean, by content:
%    III §3  "(se III, 2)"  dynamic/symbolic truth ....... = II §2   emended
%    III §4  "(II, 3)"      spirit/matter, method ........ = I  §3   emended
%    III §4  "(II, 3)"      potential psychical energy ... = I  §2/3 emended
%    III §5  "(Smlg. II, 1)" Plato's Ideas ............... = II §1   CORRECT
%    III §5  "(Smlgn. III, 4)" direction of change, time .. = II §4   emended
%    III §5  "(III, 5)"     world-process/thought-work ... ambiguous  KEPT
%    IV  §6  "(III)"        rationality of existence ..... = III      CORRECT
%    IV  §6  "(IV, 4)"      metaphysical idealism ........ = III §4  emended
%    IV  §6  "(III, 1)"     the three kinds of understanding = II §1  emended
%  (All nine verified against the 1902 print, so none of them is an e-book
%  artefact; they are Høffding's own.)
%  The offset is not systematic.  "(Smlg. II, 1)" is right as printed and
%  cannot be reconciled with "(III, 1)", which points at the same section; and
%  the bare "(III)" in IV §6 is right while the two chapter-and-section
%  references beside it in the SAME paragraph are each one chapter high.  No
%  rule about Høffding's numbering will sort them: only content will.
%  RULE ADOPTED: emend a reference to the section it demonstrably means, proved
%  by content, correcting the chapter numeral and leaving the section numeral
%  alone; where the target cannot be established, keep the reference as printed.
%  Hence "(III, 5)" stands: it fits II §5 under the offset, but also reads as a
%  backward glance at the section it sits in, and nothing decides between them.
%  All references above appear in the English in the book's OWN printed chapter
%  numbering, so they can be followed.
%
%  THE BRACE IN II §5 (p. 51).  Høffding's series of subject- and
%  object-determinations is printed "S1 { O1 { S2 { O2 ....", with a LEFT BRACE
%  between the terms.  Verified on the 1902 print itself (scan p. 51, read off
%  the page image at 400 dpi, not from OCR).  Set as printed.
%    This corrects an emendation made earlier from Høffding's later books, when
%  no scan was to hand.  Den menneskelige Tanke (1910) §66 defines a sign "<"
%  for "ethvert inkonvertibelt, transitivt og asymmetrisk Forhold" and quotes
%  this very series back as "Rækken S1 < O { S2 …", mixing the two characters
%  inside one formula; that looked like proof the brace was an OCR error for
%  "<".  It was not.  The 1902 book uses the brace; the "<" is Høffding's LATER
%  notation, and the 1910 quotation is either his own modernisation or its
%  scan's OCR.  Two independent OCR passes — the SAGA e-book and the KB scan —
%  both read "{", and the page image agrees with them.
%    The lesson is worth keeping: agreement between Høffding's later usage and
%  a plausible OCR-confusion story was not evidence about what the 1902
%  compositor set.  Check the page.
%  Subscripts: the e-book flattens "(S0)" and "(Os)" to inline letters; the
%  print sets them as subscripts, lower case ("hverken et So eller Os kunde
%  bestaa").  Set here as \textsubscript, lower case.
%
%  FISHER 1905.  An authorized English translation exists: Galen M. Fisher,
%  preface by William James (Macmillan, 1905), archive.org/details/
%  in.ernet.dli.2015.264599.  It is consulted for TERMINOLOGY ONLY.  It is not
%  a model for structure: Fisher compresses freely, and in the Introduction
%  alone he drops whole clauses.  This translation renders the Danish in full.
% ============================================================
\documentclass[12pt,a4paper]{book}
\usepackage[utf8]{inputenc}
\usepackage[T1]{fontenc}
\usepackage{libertinus}
\usepackage{libertinust1math}
\usepackage{textalpha}    % occasional Greek carried verbatim from the Danish
\usepackage[english]{babel}
\usepackage[protrusion=true,expansion=true]{microtype}
\usepackage{setspace}
\usepackage[margin=1.2in]{geometry}
\usepackage{fancyhdr}
\setlength{\headheight}{15pt}
\usepackage{enumitem}
\usepackage{hyperref}

\hypersetup{
  pdftitle={Philosophical Problems},
  pdfauthor={Harald Høffding},
  hidelinks
}

\pagestyle{fancy}
\fancyhf{}
\fancyhead[LE,RO]{\thepage}
\fancyhead[RE]{\textit{Philosophical Problems}}
\fancyhead[LO]{\textit{\leftmark}}

\setstretch{1.3}

\title{\textbf{Philosophical Problems}}
\author{Harald Høffding\\[6pt]
  \normalsize translated from \textit{Filosofiske Problemer} (1902)}
\date{København: Det Nordiske Forlag, 1902}

% --- helper: unnumbered chapter that still enters the ToC and running head
\newcommand{\fpchapter}[1]{%
  \chapter*{#1}%
  \addcontentsline{toc}{chapter}{#1}%
  \markboth{#1}{#1}%
}

\begin{document}
\maketitle
\thispagestyle{empty}
\newpage

\tableofcontents
\newpage

% ============================================================
%  INTRODUCTION  (pp. 1--5; notes 1--3)
% ============================================================
\fpchapter{Introduction}
\label{ch:indledning}

Philosophical ideas can, as the history of philosophy shows, have a twofold
significance. They can be attempts at the formulation, the treatment, or the
solution of certain problems, and they can stand as symptoms of certain
tendencies in the intellectual life of man. There is a constant interaction
between these two sides of philosophy, in that it may stand as a symptom of a
distinctive intellectual movement that a problem is formulated and treated in a
particular way, while on the other hand the sting of the problems can release
intellectual movements that would otherwise never come to be. There is expressed
in this interaction an inner connection between personal life and scientific
inquiry, one which asserts itself more or less in every scientific field, but
which in the nature of the case will come forward quite particularly as we
approach the border regions of the sciences. It comes forward more in the human
sciences than in natural science, and it comes forward most of all in
philosophy, whose problems are essentially border problems.

It is admittedly no uncommon assumption that there is a conflict between
personality and scientific inquiry --- that one had best divest oneself of one's
personality when one wishes to think scientifically, and that on the other hand
scientific methods and results signify nothing for personal life in its free
development. I myself in my youth paid homage to such an assumption, under the
overwhelming influence of S.~Kierkegaard's ideas. That I have come more and more
away from it is, I venture to say, because I have learnt to know both science
and personal life better.

The need for scientific inquiry is a special form of the need for accord with
oneself under all the manifold and shifting experiences. It expresses itself both
in the formal sciences and in the real --- both as a need to form series of
thoughts in which the one member develops with inner necessity out of the other,
and as a need to connect the experiences before us in such a way that as
comprehensive and as firm a continuity as possible is attained. Personality
consists first and foremost in an inner unity and connection among all
representations, feelings, and strivings. It is then no sacrifice that
personality makes when it stakes its life upon inquiry. It brings forth in
science --- as in art --- an objective counter-image of itself. Only thus is
the inward character that the joy of knowledge can take on to be understood, the
\textit{amor intellectualis} that the inquirer feels in working his way through
to a standpoint from which the particular can be seen as a member of a great
connection. The sacred fire which ever anew sets thought in motion, in spite of
all that can damp and hinder it, is understood only when one has this in view.
And personality on its side needs to be purified by bowing before the objective
connection of thought, the fixed order of things, within which every single
being has its determinate place. Freedom is won through truth. It belongs to the
nobility of personality that it can by its own power find the great laws which
condition its subsistence, determine the wishes that arise in it, and determine
the conditions under which those wishes can be realized. There is a science of
personal life itself, as there is of everything else; this science is perhaps
more imperfect than other science, but it is presupposed in every serious
discussion of the forms and demands of personal life. Even Charles Renouvier,
the philosopher now living who lays the greatest weight on the opposition
between necessity of thought and personality, nevertheless requires ``a rational
concept of personality''\footnote{Ch.~Renouvier: \textit{Les dilemmes de la
métaphysique pure.} Paris 1901, p.~202.} and thus cannot reject every analogy
between them.

Such an analogy will come clearly forward when one attempts to formulate the
principal problems with which philosophical inquiry occupies itself. These
problems are posed at one and the same time from the side of personality and
from the side of science.

By a purely historical route I have sought, in my work \textit{Den nyere
Filosofis Historie}, to show that there are four such principal problems, namely
I) the problem of the nature of conscious life (the psychological problem),
II) the problem of the validity of knowledge (the logical problem), III) the
problem of the nature of existence (the cosmological problem), and IV) the
problem of valuation (the ethical-religious problem).\footnote{These four
problems contain the ancient tripartition into Logic, Physics [i.e.\ Cosmology]
and Ethics, deriving from Plato's school (according to \textit{Sextus
Empiricus}: \textit{Adv.\ mathematicos} VII, 16, from Xenocrates), with the
addition of the psychological problem, which in more recent times has pushed
itself forward as an independent point of departure.} It will now be my task in
this treatise to show the inner connection among these problems. It is one and
the same fundamental problem that appears in them in different forms and
applications.

It is very different motives that can lead to philosophical inquiry. --- Very
often it is an ethical-religious, and so a practically personal, interest that
leads to philosophy. One then begins with the fourth problem. But it soon shows
itself that the treatment of that problem requires insight into the nature of
conscious life, into the conditions of knowledge, and into the constitution of
the existence in which the valuing personality exists. --- If it is an interest
in observation that lies at the ground, one will begin with the psychological
problem, as has happened in empirical philosophy in its various forms. --- If on
the other hand it is a need to distinguish between what can be known and what
cannot be known, one will begin with the problem of knowledge; this is the road
critical philosophy takes. --- If one has full confidence in the power of
thought and seeks a completed world-view, one will --- like the dogmatic and
speculative tendencies --- begin with the problem of existence. --- It will of
course not be a matter of indifference, for the kind and manner of the treatment
of the problems, which problem one begins with; the one problem will easily be
put in the shade by the other. For a comparative doctrine of problems, such as I
here seek to contribute to, what matters is to let each single problem come into
its full right. I proceed, then, in this way: in keeping with the analogy
indicated above between personality and science, I first bring forward the
problem of the nature of personal life, that is, the psychological problem, and
then pass over to investigating the nature of scientific knowledge, that is, the
problem of knowledge. When the problem of existence is placed third in the
series, the transition to it is natural both from personality, which stands as a
single part of the whole of existence, and from science, whose task it is to
lead to a world-view. The three problems named so far would be posed even if man
were a purely intellectual, merely knowing being. The fourth problem is posed on
account of the relation in which man, as a feeling and willing being, stands to
existence. Thus the principal problems of thought hang together with man's
theoretical and practical interests.

But more important than the question of the order in which the problems are
presented is the question whether they can be led back to one pervading
fundamental problem. The possibility of this I find in the significance which
the relation between continuity and discontinuity has within each of the several
problems. This relation expresses the deepest interest both of personality and
of science. In both domains there is a striving, as already remarked, after
unity and connection, and the discontinuous stands so far as an obstacle to be
overcome. But on the other hand it is discontinuity (difference in time, in
degree and in place, difference of quality, difference of individuality) that
everywhere, in the domain of science as in that of life, brings new content,
releases the bound forces, and sets the great tasks. Neither of the two elements
stands, then, as from the outset the only one justified. It will have its
unquestionable interest to follow the relation between them under the four
different points of view that our four principal problems indicate.

In the philosophy of the nineteenth century the significance of the problem of
continuity comes forward in a characteristic way, in that the different
tendencies plead the cause of continuity by turns. In the first half of the
century philosophical idealism asserted in its fashion the continuity in
existence and looked down on the science of experience on account of the
latter's fragmentary character, while positivism (as Comte and Stuart Mill
maintained it) impressed the discontinuity between the different groups of
phenomena. At the century's close it is realism that, with the help of the
hypothesis of evolution, would assert continuity, while the idealist tendency is
inclined to lay weight on the unavoidable discontinuity in our
knowledge.\footnote{Compare, besides the work of Renouvier's named above
(Note~1), the interesting writings of \textit{Emile Boutroux}: \textit{De la
contingence des lois de la nature.} 2.~éd. Paris 1895, and \textit{De l'idée de
loi naturelle dans la science et la philosophie contemporaines.} Paris 1895. ---
Related is the standpoint of \textit{Rudolph Eucken} (\textit{Der Kampf um einen
geistigen Lebensinhalt.} Leipzig 1896).} Thus the tendencies change places in
the great single combat in which the fight for truth is to be waged. When a
significant thought can no longer be carried through from one point of view,
inquiry involuntarily seeks a new point of view, and thought gains in clarity
through these shifting points of view. The different tendencies relieve one
another, like the runners at the ancient torch festival, but the torch is the
same. And if none of the tendencies that have so far come forward should have
done full justice to the fundamental thought in all the philosophical problems,
there is so much the more reason to work at forming it with greater clarity.

\begin{center}---\end{center}

% ============================================================
%  I.  BEVIDSTHEDSPROBLEMET  (pp. 6--27; notes 4--25)
% ============================================================
\fpchapter{I.\quad The Problem of Consciousness}
\label{ch:bevidsthed}

\noindent\textbf{1.}~When we begin the discussion of the philosophical problems
by investigating the concept of personality itself, we begin with psychology,
and it is then first of all a matter of determining psychology's relation to
philosophy. We can here no doubt leave aside the view, which \textsc{Herbart}
and \textsc{Lotze} still maintained, that psychology should be dependent on the
view one arrives at concerning the problem of existence, and so dependent on
``metaphysics'' or cosmology. On all sides an effort is at work to assert the
independence of psychology. But even where this effort prevails, there have
still been extremely different views about psychology's position.

There are those who hold that psychology is one with philosophy as a whole, and
that epistemology, metaphysics, aesthetics and ethics are its different parts.
Such a standpoint was taken in earlier times by \textsc{Fries} and
\textsc{Beneke}, and is taken in our own day by \textsc{Th.~Lipps}.\footnote{%
\textsc{Th.~Lipps}: \textit{Psychologie, Wissenschaft und Leben.} München 1901.
On Fries and Beneke see my work \textit{Den nyere Filosofis Historie} II,
p.~221\,f.; 236; 240.} This standpoint is untenable, partly because personal
life is only one of the subjects that can be an object of philosophizing, partly
because psychology itself, like all other science, presupposes the general forms
and principles of knowledge. The problem of consciousness thus necessarily
points beyond itself to the problem of existence and the problem of knowledge,
and for all its independence psychology is only one discipline among several.
Psychology can then only be a part of philosophy, not constitute the whole of
philosophy. It is legitimate to begin with psychology, since psychology
describes the place, and the presuppositions, from which we seek to orient
ourselves in existence.

But does psychology belong to philosophy at all? Is it not a special science
which in and for itself has no more to do with philosophy than natural science
or history has? The special roads along which psychological knowledge is sought
in the most recent period --- the experimental, the physiological and the
historical method --- do after all seem to testify that psychology is now in the
course of becoming a special science and must be separated off from philosophy!
--- To this the answer is, so far as I can see, that psychology, for all its
special methods, always presupposes the capacity for self-observation. It is
subjective, descriptive and analysing psychology that sets the tasks which can
then be treated in detail with the help of the special methods. There is here a
foundation common to psychology as a special science and to philosophy.
Philosophy must build on a general concept of the nature of conscious life if it
is to be able to treat the problems of knowledge, of existence and of valuation.
It is personality's position as knowing and valuing, and as a part of the whole
of existence, that these problems concern, and a concept of personality is
presupposed which none of those special methods is able to furnish. They
investigate single expressions of conscious life, but their investigations must
be worked together if a concept of personality is to be formed.

But here precisely the fundamental problem arises in the psychological domain:
is it possible by the road of experience to arrive at a concept of personality?
Does our conscious life form a whole, a continuum, a little world of its own, or
should it be merely an aggregate, a sum of elements and fragments? This is the
question that philosophical psychology (or, if one prefers, the philosophy of
psychology) raises. And in its treatment of it it relies partly on
self-observation and descriptive psychology, partly on the results of the
experimental, the physiological and the historical method.

Experience shows us greater and lesser interruptions of our conscious life;
there are unconscious intervals between the conscious states. And when we
investigate the single state of consciousness more closely, we can distinguish
within it different elements which can be found again in other states, so that
the single state, no less than consciousness as a whole, might seem to be a
resultant of a manifold of elements or fragments which stand as what is properly
real. Consciousness would then be an aggregate, or a sum of single flashes. A
psychical continuum, such as is presupposed in the concept of personality, there
would not be.

To this it must first be remarked that it will never be possible to decide with
full certainty whether analysis has reached down to the last elements. The
possibility always stands open that the elements at which we stop, and out of
which we are inclined to think consciousness composed, are themselves composite,
and that if we could reach down to still simpler elements, elements of the second
or third order, a continuity would appear which could not be exhibited so long
as we stopped at the elements of the first order. Our sensations, which one has
been inclined to regard as a kind of psychical atoms, nevertheless bear --- as
recent investigations along various roads have shown --- traces of processes of
composition of a still more elementary kind than those which observation shows
us more directly.\footnote{See my \textit{Psykologi} V~A,~2.}

This, however, is only a preliminary consideration, and leads to no settled
result, since the defender of discontinuity may very well use it to show that
all the continuity that can be exhibited is only preliminary. What is on the
other hand of decisive importance is that the so-called psychical elements are
always determined by the connection in which they occur, and that it is only by
abstraction that one ascribes to them, outside this connection, properties which
they have only within the connection itself. It is on this fundamental thought
that my presentation of psychology is built. All mental life works in the manner
of genius: it brings forth immediately and involuntarily connected phenomena and
processes which analysis afterwards seeks, in a more or less artificial way, to
dissolve into ``elements''. As regards sensation we can here appeal to Fechner's
law and to the law of contrast, according to which the degree of intensity and
the quality of every single sensation are determined by the whole connection in
which it comes to be. The connection cannot, then, be conceived as a product of
the psychical elements, since these come to be as \emph{these} determinate
sensations only in virtue of the connection. Did they occur in another
connection, they would not be the same sensations. The relation is analogous as
regards representations. The association of representations finds its ultimate
explanation in the unnaturalness of isolating the single representations; there
is always at work --- and to a higher degree, the more vigorous consciousness is
--- a tendency to supplementation or extension, whereby the single
representation enters into combination with other representations according to
determinate laws. The so-called ``Association-Psychology'' (among whose adherents
I have myself sometimes wrongly been reckoned) conceived the single
representations as independent atoms which were brought into connection with one
another in a purely external, mechanical way. But it is precisely the reverse:
it is the connection of the whole that in the process of association exercises
its power over the single representations, which moreover never occur in a
wholly isolated condition, such that their combination could stand as a
mechanical product. Analysis here too always presupposes synthesis. In thinking
proper this relation shows itself plainly when we compare the formation of
judgments and the formation of concepts. These mutually presuppose one another,
since our judgments, which are after all combinations of concepts, can only be
complete when the concepts combined in them are complete, while conversely the
formation of concepts presupposes a series of judgments by which different
elements are determined in their mutual relation. It shows itself here again
that the whole and the parts mutually determine one another. There are no wholly
isolated concepts which could only afterwards be combined into judgments.

There expresses itself here an antinomy which hangs closely together with the
essence of consciousness and is peculiar to the concept of personality.
Consciousness and personality cannot be explained as products of antecedently
given elements, any more than organic life can be explained as a product of
inorganic elements. But on the other hand the essence of consciousness and of
personality, like the essence of organic life, displays itself in a constant
gathering-together of \emph{given} elements, elements which they do not
themselves originally bring forth. It is this antinomy that makes the
coming-into-being of life and of personality such great riddles.\footnote{%
Compare my treatise \textit{Det psykologiske Grundlag for logiske Domme}
(Vid.~Selsk.~Skr.\ 6.~Række), p.~16\,f.; 59, and my treatise \textit{Om
Vitalisme} (reprinted in ``Mindre Arbejder'').}

In the foregoing, account has been taken only of the more formal connection
within conscious life. But in every consciousness there is a goal towards which
there is striving, a guiding interest which takes, or strives to take,
everything into its service. This guiding interest --- the fundamental value,
one might call it --- can change in the different ages of life; but there will
always anew be a tendency to the development of such a one, and in higher or
lower degree it sets its stamp on all the elements of consciousness, gives them
their direction. It constitutes the soul proper (when by ``spirit'' we think
chiefly of the formal connection).

It is the relations here brought forward that make the concept of personality
have always to stand as the central thought in psychology. When analysis and the
special methods make their sections and seek to isolate single elements or
moments, that is legitimate, just as it is legitimate to attempt to determine an
irrational number by adding decimal to decimal, as far as one can. But the
irrational relation between the whole and the elements continues to subsist.

From the idealist side one has often been inclined to take the concept of
personality for something finished and to operate with it in cosmological
speculations. What is thereby overlooked --- and it is what the positivist
tendency particularly emphasizes --- is that we are constantly at work forming
the concept of personality, spelling our way forward to a psychological
conception of the whole, just as biology is still constantly spelling out a
definition of life. And just as biology, for all its recognition of the
peculiarity of the living organism, knows no other method than that which
through observation, experiment and analysis seeks to understand the composite
processes by the help of the simpler, so psychology, however firmly it holds to
the synthetic character of conscious life, can apply no other methods than those
which all science applies: observation, experiment and analysis. The concept of
personality stands as the ideal towards which the course is set, the standing
problem which all the special methods serve to illuminate.\footnote{Of interest
for the standpoint of modern positivism on this problem, in contrast to the
standpoint Comte and Stuart Mill took, is \textsc{Roberto Ardigò}'s work
\textit{L'unità della coscienza.} Padova 1898.}

The irrational is here, as everywhere, not merely a barrier but also a task, a
task that is set ever anew. Descriptive psychology, which leads one particularly
to lay weight on the connection of the whole in and with which mental phenomena
appear, will here --- especially as regards the higher or more developed
phenomena of consciousness --- always preserve its independence and its
importance in relation to experimental psychology; the latter in the end always
receives its tasks from the former. And descriptive psychology serves at the
same time as a corrective to experimental psychology, which by its method
easily comes to isolate single elements too much, to overlook what is purely
involuntary in conscious life, and to lay too great weight on the outward
symptoms of states.\footnote{Compare my \textit{Psykologi} I,~8~d.} On the other
hand descriptive psychology will be able to draw advantage from experimental
psychology, in that what the latter is able to exhibit concerning the more
elementary psychological processes can by the road of analogy throw light on the
nature of the higher processes and give description greater fullness and
accuracy.

Descriptive psychology approaches art and becomes in the end an art;
experimental psychology approaches natural science. But there is nevertheless no
difference of principle between them. There is not, as \textsc{Münsterberg} in
the most recent period especially has wished to maintain,\footnote{%
\textsc{Münsterberg}: \textit{Psychology and Life.} Westminster 1899. ---
Münsterberg allies himself with the view that Windelband and Rickert have
urged, according to which there is a gulf between natural science (to which
psychology too is reckoned) and the science of culture. On this view see the
apt remarks in \textsc{Wundt}: \textit{Einleitung in die Philosophie.} Leipzig
1901, p.~65--74, and in \textsc{Guido Villa}: \textit{Psychology and History.}
(The Monist, January 1902.)} a gulf between them. The former is to be a
``science of value'', the latter a natural science; the former is to operate
with the ``concept of freedom'', the latter with the concept of cause. By
freedom Münsterberg understands the capacity to act according to purposes. But
should then the psychical process by which a human being comes to set himself a
purpose and to work at attaining it not be an object for scientific psychology?
And can it be understood otherwise than by being investigated in connection with
the simpler processes that take place when pleasure and pain, joy and sorrow
arise, or when representations are combined and separated? Here there is a mass
of experience to be gathered and collected; what matters is to follow a
development step by step, to establish the different degrees of consciousness
and the differing impulses and motives that co-operate. There will further be a
distinction to be drawn between the different individual types that can appear
with respect to the kind and direction of the life of will, and a comparative
individual psychology will be able to serve as a necessary aid to understanding
the development within the consciousness of the single person. It is then a
wholly unnatural boundary that is drawn when one lays down a difference of
principle between ``personal'' and ``psychophysical'' categories --- and
especially when this opposition is made the ground of an opposition between
``the truth of life'' and ``the truth of science''.\footnote{%
\textsc{Münsterberg}, op.\ cit.\ p.~282: ``If we force the system of science
upon the real life, claiming that our life is really a psychophysical
phenomenon, we are under the illusion of psychologism. If, on the other hand, we
force the views of the real life, the personal categories, upon the scientific
psychophysical phenomena, we are under the illusion of mysticism. The result on
both are the same. We lose the truth of life and the truth of science.''} It is
after all the task of science to understand life, even if complete understanding
is always an ideal, and it is life that yields science the material of
experience on which it builds. Here there is a constant relation of interaction,
but no absolute opposition. And there are many intermediate links and
intermediate forms between science's separating out of the single elements of
existence and life's great interplay of all elements, which in the end only art
can present in its whole fullness.

\noindent\textbf{2.}~Nevertheless great difficulties always remain for the
scientific treatment of the phenomena of consciousness. To a real understanding
it belongs that a continuity be exhibited in the processes one investigates. It
would be psychology's ideal not only that we could win so complete a description
of the states of consciousness that each single one of them came to stand as a
distinctive member of the whole psychical process, but also that we could reduce
the differences between the changing states to forms so simple that the
following appeared as continuations, or as transformations, of the preceding.
But the prospect of reaching this ideal is barred by the discontinuity which
experience seems to establish. There are, after all, unconscious intervals
between our conscious states --- in fainting, in dreamless sleep (so far as
there is such a thing); and there are qualitative differences between the
different states and elements of consciousness, so that every single state and
every single element seems to come to be out of nothing, if we keep to the
strictly psychological consideration. Between the different individual
consciousnesses there is in any case a great and conspicuous discontinuity. The
one consciousness can be derived from the other still less than the one state of
consciousness can be derived from the other.

These psychical relations of discontinuity come forward especially sharply when
they are set over against the continuity, and the relations of equivalence,
which the material phenomena present and of which men became aware very early.
Among the Greek natural philosophers the persistence of matter through all
changes was a common assumption, and this assumption has received its
confirmation in the chemistry of more recent times. In sharp contrast to it
stood the distance between souls among themselves. When \textsc{Descartes}
applied the concept of substance both in the psychical and in the physical
domain, it was with the great difference that he assumed many soul-substances
but only one single material substance.\footnote{\textsc{Descartes}
(\textit{Synopsis meditationum}) teaches that all substances are imperishable,
but that it is only \textit{corpus in genere sumptum} that is substance, not for
example \textit{corpus humanum}; every soul, by contrast, is \textit{pura
substantia}, which does not change, however much its single \textit{accidentia}
(thoughts, feelings and expressions of will) may vary. --- In his work
\textit{Cogitata metaphysica} (II c.~11) \textsc{Spinoza} says (from the
Cartesian standpoint which in that work he essentially adopts): \textit{Si ad
totam Naturam materiæ attendamus, illi nihil novi accedit; at respectu rerum
particularium aliqvo modo potest dici, illi aliquid novi accedere. Qvod an etiam
locum habet in rebus spiritualibus, non videtur: nam illa} [i.e.\
\textit{spiritualia}] \textit{ab invicem ita dependere non apparet} [souls do not
depend upon one another as bodies depend upon one another]. Later Spinoza came
to another view.} Thereby the psychical discontinuity was pronounced sharply and
definitely, in contrast to the physical continuity. It is only in a half or
wholly mystical fashion that we find, in the history of philosophy, the
continuity of mental life asserted in such a way that the single consciousnesses
are thought to stand in the same relation to a universal soul-substance in which
the single bodies stand to the universal material substance. So with
\textsc{Averroes}, \textsc{Spinoza} and \textsc{Hegel}.\footnote{Relying on his
interpretation of Aristotle's ``active reason'', \textsc{Averroes} (on him see
\textsc{Renan}: \textit{Averroës et l'Averroïsme.} Paris 1852) teaches that while
the single souls arise and perish, there subsists \textit{intellectus
universalis}, the world-thought that is at work in the thinking of the single
souls. In the Middle Ages this doctrine was renewed, among others, by
\textsc{Siger of Brabant} in his \textit{Qvæstiones de anima intellectiva},
recently edited by Mandonnet. (Mandonnet: \textit{Siger de Brabant et
l'Averroïsme latin au 13\textsuperscript{e} siècle.} Fribourg 1899.) The whole
doctrine of the actual subsistence of spirit goes back to Neoplatonism. (See
\textsc{Plotinos}: \textit{Ennead.} V, 9, 6.) --- On \textsc{Spinoza}'s doctrine
of \textit{intellectus infinitus} and \textit{idea dei} see \textit{Den nyere
Filosofis Historie} I, p.~290; on \textsc{Hegel} ib.\ II, p.~165.} In the most
recent period it is now not so much the relation of discontinuity between
individual consciousnesses as the relation of discontinuity within the single
consciousness that draws attention to itself. And there are in particular two
considerations which have here led to the problem's being posed more sharply
than before.

Physiology has come to full consciousness of its method and of its independence
as a science, especially after the reform it underwent about the middle of the
nineteenth century. It now consistently requires that every state of the brain
be explained from the preceding states of the brain and of the organism ---
that physiological phenomena in general be explained by physiological causes, in
so far as they are not effects from the physical world outside, continued into
the organism. On strict natural-scientific method every material state within
and outside the organism must find its explanation by being seen as having
arisen through the conversion, into a new form, of the energy which expressed
itself in the preceding states. When two brain-states succeed one another, the
explanation will accordingly first have been reached when it is shown that all
the energy which expressed itself in the first state has been converted without
remainder into nutritive activity, heat, electricity, motor impulses and so on.
It is indeed difficult to imagine such an experiment carried out; but it would
mark the carrying through of the physiological method. Neo-vitalism itself knows
no other method for the investigation of physiological states. Thus is confirmed
\textsc{Spinoza}'s saying, that when one appeals to the intervention of the soul
in explanation of a bodily state, it is the same as saying that one does not
know what the cause is. Even if an experiment should show that at the arising of
an organic state energy arose or disappeared without physical equivalent, one
would surely sooner believe that there was a fault in the experiment than
acquiesce in the result. ``The principle of the conservation of energy has'',
says \textsc{Maxwell},\footnote{\textsc{Maxwell}: \textit{Scientific papers.}
Vol.~II. Cambridge 1890, p.~759. Shortly before the passage quoted, Maxwell
remarks that whereas the thought that the soul might be found somewhere in the
brain by anatomical means has nowadays as a rule been given up, the thought has
held on longer that by pursuing the material processes far enough back one will
come to a material process brought about by the soul. It is against this
possibility that the quoted passage is directed. Compare \textsc{Spinoza}:
\textit{Ethica} III, 2, Schol. --- His conception of the principle of inertia in
his work \textit{Matter and Motion} §~41 leads to a result similar to Maxwell's
conception of the principle of energy. --- Compare my
\textit{Psykologi}\textsuperscript{4} p.~30 and my \textit{Psykologiske
Undersøgelser}, p.~82--86.} ``acquired so great a scientific weight that no
physiologist would have any confidence in an experiment which showed a
considerable difference between the work performed by a living being and the
reckoning of energy received and expended.'' --- By carrying through the
physiological method, understood in this spirit, an ever greater continuity will
be exhibited in the organic processes to which the phenomena of consciousness
in fact show themselves to be attached. Physiology is therefore far more happily
placed with respect to the principle of continuity than psychology will ever be
able to be, and it might therefore seem a well-founded demand that \textsc{Carl
Lange} put forward when he declared that all psychological definitions ought to
be replaced by physiological ones.\footnote{\textsc{Carl Lange}:
\textit{Nydelsernes Fysiologi.} København 1899, p.~45.} The psychical phenomena
then come to stand as mere provisional symptoms.

The second consideration attaches to the one just adduced, but is of a general
epistemological kind. It rests on the proposition that a complete understanding
is won only where the relation between cause and effect can be made a relation
of identity or of continuity, such that a quantitative equation becomes possible
between the phenomena that are connected in the causal relation. It is the ideal
concept of cause that is here maintained, in contrast to the empirical concept
of cause, by which phenomena that are qualitatively different stand as connected
in unavoidable succession. The closer investigation of these two concepts of
cause I postpone to the problem of knowledge. Here it need only be remarked that
there is no doubt that where the ideal concept of cause can be carried through
we obtain a more complete understanding than where we must stop at the empirical
concept of cause --- and have, perhaps, difficulty enough even in applying that.
Now both the material and the psychical phenomena present qualitative
differences; but in the material domain it is possible to conceive these
qualitative differences as differences of quantity, namely by virtue of the law
of the conservation of energy and of matter, while nothing corresponding is
possible in the psychical domain. The right psychological method is therefore
taken to consist in substituting the physiological phenomena for the
corresponding psychical ones; that is the only way in which we can put
quantitative determinations in the place of qualitative ones and thereby carry
through a stringent causal connection. This conception of the psychological
problem has been carried through most strictly by \textsc{Richard Avenarius} in
his acute work \textit{Kritik der reinen Erfahrung}. He requires there, as the
condition of full scientific understanding, a reduction of ``the dependent vital
series'', by which are meant the psychical states, to ``the independent vital
series'', the corresponding physiological states, whereby alone it becomes
possible to reduce the empirical differences to the least possible, or, as
Avenarius expresses it, to reach ``a heterotic minimum''. In a similar direction
goes \textsc{Münsterberg}'s view, as he has expressed it in his work
\textit{Psychology and Life}. His principal contention is that since psychical
phenomena are not quantitative, they cannot be members of a causal
series.\footnote{\textsc{Münsterberg}: \textit{Psychology and Life.} p.~127:
``Mental facts, as they are not quantitative, cannot enter into any causal
equation.''}

As against the views which thus require a reduction of psychology to physiology
in order that a scientific psychology should become possible --- which really,
then, would abolish psychology in order to make it a science --- it must be
acknowledged that they can appeal to the factual discontinuity and the
qualitative difference between psychical phenomena, which will always more or
less set the limit to a strict carrying through of psychology as a science.
Whether the road they point out leads to the goal, and whether the method they
commend does not contain a palpable self-contradiction, is another matter.

When it is required that the psychological definitions be \emph{replaced} by
physiological ones, it is clear that the presupposition of this is that the
psychological definitions must first be finished. But to procure those
definitions must be psychology's business, and if it cannot procure definite and
secure definitions, physiology cannot know what it is that it is to seek an
explanation of in the brain: if what is to be replaced is vague and wavering,
the replacement will be so accordingly. There will be no security that one has
really found the brain-states which correspond to the psychical phenomena one
has in view. Psychology's independence must therefore in any case be
acknowledged, since it is psychology that here --- as a kind of symptomatology
--- is to set physiology its tasks. Now for every science of experience it is a
long and difficult labour to find conclusive definitions; such are possible only
when the science of experience is in essentials completed; they come at the end,
not at the beginning, of the investigation. Far too often one operates with
unfinished psychological definitions which are regarded as trustworthy points of
departure for the physiology of the brain. Thus \textsc{Descartes}, and in more
recent times \textsc{Lotze}, regarded the very concept ``soul'' as so distinct
and clear that one could confidently invite physiologists and anatomists to
search for ``the seat of the soul''. In the most recent period
\textsc{Flechsig}\footnote{\textsc{Flechsig}: \textit{Gehirn und Seele} 2.
Leipzig 1896, p.~24\,f.} has regarded the concept ``association'' as so simple
and clear, and withal as so independent, that he has been able to think he could
find a peculiar place in the brain for the functions that concept designates.
That it is nevertheless a highly imperfect concept of association that Flechsig
here takes as his ground is seen from the fact that psychological experience
never shows us absolutely single psychical elements which are only afterwards to
be brought into connection with one another by a peculiar process. One cannot
set the process of association in definite contrast to the processes by which
the single psychical elements (sensations and representations) come to be. It is
a psychological abstraction, not a real psychological definition, that Flechsig
operates with. Anatomical objections have been raised from several sides against
Flechsig's doctrine of peculiar association-centres, which I cannot judge; but
the psychologically unfinished character of his concept of association shows
sufficiently how difficult a task it is to ``replace'' psychological definitions
by physiological ones.\footnote{Compare the criticism of Flechsig's theory in
\textsc{O.~Vogt} (\textit{Flechsigs Associationscentrenlehre, ihre Anhänger und
Gegner.} Zeitschrift für Hypnotismus. V, 6) and \textsc{Alb.~Adamkiewicz}
(\textit{Die Grosshirnrinde als Organ der Seele.} Wiesbaden 1902, p.~75\,f.).}
A principal point for the criticism of Flechsig's theory seems to me to be the
following circumstance. When there are psychical elements which do not enter
into such an association as might under usual conditions be expected, then on
closer investigation the cause will show itself to be that the elements in
question are each of them members of other, fixed associative relations from
which they cannot be released. It will in any case always, psychologically
regarded, stand as a task to find the preceding associations which hinder later
associations that are in and for themselves ``natural''. By isolating the two
concepts element and association from one another one will thus get both a false
psychology and a false physiology. For an exhaustive definition of a concept
like association long investigations will still be needed, and physiology will
have to wait a long time if it means to ``replace'' a psychological reality and
not a mere abstraction.

But let us suppose that the psychological definitions are finished and complete
--- will not the greatest problems still remain unsolved? How can physiological
states have psychical symptoms? How can qualitative differences correspond to
quantitative ones, and how can discontinuous phenomena be attached to continuous
processes? Far from such questions falling away with the reduction of psychology
to physiology, they are precisely sharpened by it and appear in a more glaring
form. And in any case it stands as an unsolved riddle how the qualitative
differences and the discontinuity arise. Even if one succeeds in showing that
two concepts \textit{A} and \textit{B}, which one has hitherto regarded as
different, are identical, so that one can set \textit{A} = \textit{B}, one has
not thereby explained how \textit{A} and \textit{B} could stand as different in
the first place. To call this difference merely ``subjective'' does not help;
it is not thereby got out of existence --- it is precisely the fact that
existence has a subjective side as well that makes psychology possible and
necessary. And if one declares (with \textsc{Münsterberg}) that these subjective
phenomena can neither be described nor explained,\footnote{\textsc{Münsterberg}:
\textit{Psychology and Life.} p.~172: ``The reality of the will and feeling and
judgment do not belong to the describable world.'' --- p.~208: ``The subjective
attitude is never object; it is never perceived.''} it is not easy to understand
how one is to get anything to put into a ``causal equation''. One then behaves
like the man who saws through the branch he is sitting on.

It is inconsistent to deny a causal relation between psychical phenomena among
themselves while assuming a causal relation between the corresponding
physiological states. If one acknowledges a physiological lawfulness, one must
acknowledge a psychological lawfulness also. If the psychical phenomenon α
corresponds to the brain-state \textit{A}, the psychical phenomenon β to the
brain-state \textit{B}, and if there is a causal relation between \textit{A} and
\textit{B}, then a causal relation --- in the sense of unavoidable succession
--- must appear in psychological experience between α and β as well. And --- as
has been shown above --- it is in reality this causal relation between α and β
that leads one to look for a causal relation between \textit{A} and \textit{B}.
In \textsc{Avenarius}'s strict carrying through of the reduction of psychology
(the doctrine of ``the dependent vital series'') to physiology (the doctrine of
``the independent vital series'') it is likewise plain at every point that it is
from the ``dependent vital series'' that he infers the constitution of the
``independent'' one. For the most part what is constructed is only a
physiological schematism, offered as an ``intuitive'' presentation of what goes
on in the psychical processes. Neither physiology nor psychology is yet so
complete that one can get beyond such a schematism.

The psychical phenomena do not, however, present so great a discontinuity and
such purely qualitative differences as is often supposed. Exact observation
often leads to the discovery that where there seemed to be a psychically empty
space there was in reality a psychical content given, although it was not
grasped by attention or by recognition, or although it was forgotten
immediately after being experienced. And as regards the qualitative differences,
here too more exact observation will find more and finer nuances than those at
which one first stopped and which one perhaps confidently declared to be wholly
disparate. The continuity in conscious life can therefore be carried further
than is often assumed, and there can be a scientific danger in proclaiming too
strongly that difference of quality and discontinuity are the distinguishing
mark of conscious life. The inner connection in the content of consciousness is
a fact as well, and functions such as memory and comparison are quite as
significant for the characterization of conscious life as are the phenomena that
bear the stamp of discontinuity. Psychology's task therefore naturally becomes
that of exhibiting, so far as possible, the connection and combination of the
single elements, so that the whole is understood through the parts and the parts
through their relation to the whole. That this leads to an antinomy which in the
end makes the psychological problem insoluble for us, we have already seen. But
there is work enough to do before we come to that limit (which moreover cannot
be proved to be absolute). \textsc{Leibniz} has the merit of having energetically
maintained the principle of continuity in the psychical as well as in the
physical domain, especially by pointing to those psychical elements which ---
in relation to the clearly conscious states --- are vanishingly small, and which
only by being summed and combined bring forth a result accessible to unpractised
self-observation. A comparison between modern and ancient or medieval poetry
alone shows how many more mental nuances self-observation has now caught sight
of; it is hardly only because differentiation in the psychical domain is now
greater than in earlier times that this greater richness is found. Knowledge and
understanding of what is immediate and involuntary in mental life has increased
above all. Childlike and primitive psychology has a tendency to conceive
everything that goes on in the soul as clearly conscious and as resting on
deliberation and intention. In modern psychology the burden of proof --- in
practice as in theory --- has been shifted: it now falls on the one who asserts
that an action has been performed with deliberation and express decision. And
the transition between the unconscious and the conscious, and within the
conscious between the involuntary and the voluntary, itself takes place
continuously; the more practised self-observation and criticism are, the more
difficult it becomes to lay a finger on one definite point at which the boundary
is supposed to lie.

Leibniz was inclined to reduce all psychical differences to differences of
degree in respect of clearness and obscurity.\footnote{\textit{Den nyere
Filosofis Historie} I, p.~334--336.} This attempt was characteristic of the
intellectualism of the age of Enlightenment. But there are other ways than such
a reduction in which empirical psychology can conceive a continuous connection
despite the qualitative differences which the psychical processes present at
their different stages. Fusions and compositions take place at the different
levels of conscious life, from the simplest sensations to the highest states of
feeling. And displacements of motive take place, by which what at first has only
mediate value later acquires immediate value, or conversely. Finally, conscious
life in wholly integrated characters --- psychology's highest and noblest
objects --- shows such a dependence of everything that stirs in the soul upon
one purpose, one guiding thought, that we here find a causal connection quite as
firm and as intimate as any in the physical domain. Every single thought, every
single mood and every single impulse is in such a character plainly determined
by the whole to which it belongs, and which it itself helps to form.

Cases enough remain in which psychical continuity cannot be exhibited. The
question is whether we are therefore entitled to deny it. Apart from what
continued observation may establish, the great question arises here whether
psychical phenomena can come to be from material causes, when the concept of the
material is conceived in the way natural science has hitherto conceived it. If
the natural-scientific concept of the material includes no possibility of an
arising of psychical phenomena, then it will be fully legitimate to set up the
concept of potential psychical energy, in order to mark that we do not mean to
let go of the principle of continuity even where we cannot apply it in detail in
experience. It is indeed for a similar reason that the concept of potential
energy has been introduced in physics, despite the obscurity that attaches to
it. Everywhere it contains an acknowledgment of a limit at which one nevertheless
does not mean to give up the connection one has hitherto found. The mere fact
that psychical elements can be reproduced after an interval in which they have
not been in consciousness already necessitates the use of concepts like
``disposition'', ``trace'', ``possibility'' and the like, which in reality
express the same thing as is meant by ``potential energy''. That in applying
this concept in the physical domain we are more happily placed than in the
psychical\footnote{Compare my \textit{Psykologi} V~B,~6.} does not exclude the
legitimacy and the necessity of setting it up here. One can of course, in a
purely descriptive presentation, content oneself with stating the different
psychical phenomena that emerge at certain moments, often without observation
finding any psychical connection between them. The psychical phenomena then
stand as scattered flashes, in strong contrast to the continuous connection
which the corresponding physiological phenomena present. And since what appears
with the highest degree of continuity easily comes to stand for us as more real
than what is discontinuous and momentary, the physiological phenomena will
easily come to stand as what is properly real, the phenomena proper, and the
psychical phenomena will come to stand as something superfluous, an accidental
plus, or, as it has been expressed, as epiphenomena. Such an
``epiphenomenalism'' has surely never been set up as a theory proper; it is
merely an empirical registering of what is enigmatic in the appearance of the
psychical phenomena as compared with the more closed continuity in the series of
material phenomena. Nor indeed does it furnish any explanation; on the contrary
it abolishes both every psychological and every physiological explanation, if it
means to be more than a mere description. And even as description it easily
becomes precipitate, since it easily comes to stop too early at what is
discontinuous, without investigating whether there are not more transitions and
nuances to be found than have hitherto shown themselves. The concept of
potential psychical energy will here --- precisely by being really only the
expression of an unsolved task --- contain a standing summons to go further in
observation and analysis, in order if possible to discover the reality in which
the ``possibility'' in the end always consists.

\noindent\textbf{3.}~When weight is laid, from the psychological and the
physiological side alike, on preserving continuity in as high a degree as
possible, the identity hypothesis becomes the real working hypothesis for the
problem of the psychical and the physiological. Like every significant
hypothesis it is really only the expression of a method. We have in the
psychical and the physiological phenomena two series of states, or forms of
state, which empirically vary in a certain relation to one another, without its
being possible, however, to \emph{derive} the \emph{existence} of the one series
from the \emph{existence} of the other. Scientifically regarded, the task must
then be to conceive the two series each for itself as completely and as
continuously as possible, and to show which determinate members of the one
series correspond to determinate members of the other. We may rightly regard
members of the one series as symptoms of members of the other. Sometimes it is
the psychical, sometimes the physiological states that are the more accessible,
and from which we therefore take our point of departure. The close
belonging-together of the two series of states makes it impossible to refer them
to two different ``beings'' or ``things''; it will be natural to conceive them
as different forms of expression of one and the same ``being''. From the fact
that two properties cannot be derived from one another one has no right to
infer that they belong to two different things, especially when they vary in a
determinate mutual relation. Sulphur's form of state (solid, liquid and so on)
and its colour (yellow, brown and so on) vary in a determinate relation to one
another and yet cannot be derived from one another, although one can, when one
knows the two series of changes, infer from the one to the other. In the case of
sulphur we do know the common cause of the two series of changes, namely heat;
it is heat that expresses itself in these two different ways. In the relation
between the psychical and the physiological series the analogous knowledge is
denied us. Our experience is here insufficient to solve the problem. This hangs
together with the fact that both our knowledge of the psychical and our
knowledge of the physical are empirically limited. What matters, then, is to
work in both fields of knowledge in such a way that at no point does one
arbitrarily cut the thread.

The physiological continuity is both a consequence of the law of the
conservation of energy (really already of the principle of inertia, on the
strict interpretation of it) and an expression of physiology's independence as a
science. It furnishes a fruitful and indispensable principle of research by
requiring constant demonstration of physiological causes of, and physiological
effects of, physiological states. As we have seen above, the principle we stand
at here is so essential that --- if an experiment pointed in the opposite
direction --- one would sooner assume an error in the experiment than a breach
of the principle. The principle of continuity --- and with it the identity
hypothesis --- would be refuted if it could be shown that the energy contained
in a brain-state did not stand in a relation of equivalence to preceding and
succeeding states in the brain, in the organism, and in the physical world
outside. It would already be a matter of concern for the identity hypothesis if
it could be shown that the psychical phenomena in time go \emph{before} or
follow \emph{after} the corresponding physiological states. But in an attempt to
establish this it is not enough to show what popular observation already shows,
that sensation arises later than the corresponding sense-impression, or earlier
than the physiological (motor, and especially vasomotor) expressions which in
one way or another hang together with it. For between the psychical states and
the peripheral processes in the sense-organs and the organs of movement lie the
central physiological processes, which are not so readily accessible to exact
and direct experiment. And even if such a direct experiment were possible, it
would not be decisive: because the phenomenon \textit{B} follows \textit{A} in
time, \textit{A} and \textit{B} may very well be effects of the same cause,
which first deposits the effect \textit{A} and then the effect \textit{B}. There
might be peculiar causes which made the one vital expression alter more rapidly
than the other. --- Finally it would be a proof against the identity hypothesis
if it could be established that different psychical phenomena could correspond
to one and the same physiological state. It has been supposed, for example, that
the brain's significance for conscious life might consist in counteracting
organic life's tendency to habitual, automatic operation. The brain has a
wonderful capacity for initiating new movements which, as soon as they have been
practised, become habitual actions that are left to lower centres to carry out.
Only by the brain's help has thought been able to free itself from automatism
and to use the organism in its service without sinking under the power of habit.
But different thoughts can at different times make use of the same ``motor
schema'', the same suspension of an automatic tendency. This possibility has in
the most recent period been advanced by \textsc{Bergson} as an objection to
``strict parallelism''.\footnote{See the interesting discussion of \textit{Le
parallélisme psychophysique et la Métaphysique positive} in Bulletin de la
société française de Philosophie. Juin 1901. In the course of that discussion
Bergson remarked, among other things (p.~51): \textit{Étant donné un état
psychologique, la partie jouable de cet état, celle qui se traduirait par une
attitude du corps ou par des actions du corps, est représentée dans le cerveau:
le reste en est indépendant et n'a pas d'équivalent cérébral. De sorte qu'à un
même état cérébral donné peuvent correspondre bien des états psychologiques
différents, mais non pas des états quelconques. Ce sont des états psychologiques
qui ont tous en commun le même} ``schéma moteur''.} This objection touches on a
question that has come up very little in the frequent discussions of the
psychophysical problem. Whether or not one accepts the special theory the French
philosopher has set up, it seems clear that the closer one presses the problem
to life, the more extraordinarily difficult it will become to find those members
of the two series of phenomena that can be said to be ``corresponding''. Where a
qualitative difference is present, one can establish only by roundabout ways
which members correspond to one another. When we know how heat acts on a
substance's colour and on its form separately, we can determine which changes of
colour and of form ``correspond to one another''. But it is precisely an
analogous knowledge that is not at our disposal in the discussion of the
psychophysical problem. We are then limited to purely empirical inferences. And
such will here be more difficult to draw than Bergson seems to assume. For when
sufficient weight is laid on continuity in the psychological and the
physiological domain alike, it will prove difficult to resolve the two series
into members that appear with such independence and distinctness that a real
comparison can take place. But whatever hypothesis we take as our basis, we must
be prepared for it that the members of the two series which may be taken to
correspond to one another will present differences that could not be inferred
from the one member or the other \emph{taken by itself alone}. It might
therefore very well be that ``different'' psychical phenomena could correspond
to ``the same'' physiological state (or conversely), just as one language
sometimes has only one word where another has two. It is then the connection
that is decisive, and it will surely be decisive to such a degree --- in the
psychological as in the physiological domain --- that it is only with a slight
degree of approximation that one will be able to speak of ``the same'' phenomena
or states. ``Parallelism'' is therefore not touched by that objection, provided
we make sufficiently strict demands on the demonstration that the states really
are the same.

It is essentially as a working hypothesis, not as a positive solution, that the
identity hypothesis (for which ``parallelism'' and similar expressions are
unhappy and misleading designations) has its significance. I for my part have in
any case always maintained it as an ``empirical formula'' which can guide us in
our investigations in such a way that neither physiology's right nor
psychology's is violated by an investigation on the one side or the other being
broken off too early. The physiologist may be tempted to stop too early if he
can expect at some point or other to come upon ``the soul'' as the cause of a
change in the brain's state,\footnote{An example of this I have adduced in my
\textit{Psykologi}, 4.~Udg., p.~67 Note.} and the psychologist may be tempted to
stop too early if he can expect at some point or other to stand before a
psychical phenomenon supposed to have its cause in a ``nerve process'' or in
``nerve energy''.

\bigskip

\noindent\textbf{4.}~It will nevertheless have its interest, from both sides,
the psychological and the physiological, that they should approach one another
as closely as possible by going back to the fundamental concepts of the two
sides under which existence here appears in our experience.

From the psychological side the \emph{concept of will}, taken in the widest
sense as the concept of psychical activity, will stand as the fundamental
concept. It might seem to tell against this that several psychological writers,
precisely in the most recent period, seek to eliminate the concept of will from
psychology --- not because they deny what popular usage calls will, but because
they hold that this concept does not designate so fundamental a point of view as
knowledge and feeling do, the so-called phenomena of will being supposedly
reducible to elements of knowledge and of feeling in peculiar combinations. In
support of this one may appeal to the fact that will as such, our activity as
conscious beings, cannot be an object of immediate self-observation in the way
that representations and feelings can be.\footnote{Compare my \textit{Psykologi}
VII~B,~4; \textit{Psykologiske Undersøgelser}, p.~67--81; \textit{Søren
Kierkegaard som Filosof}, p.~74--81. --- \textsc{Ebbinghaus}: \textit{Grundzüge
der Psychologie} I, p.~168 (compare also his article in Zeitschrift für
Psychologie IX, p.~201), and \textsc{Ehrenfels}: \textit{Werttheorie} I,
p.~245--249, infer from the unobservability of will that will is not a special
psychical element. Compare my criticism of the last-named work in Göttinger gel.
Anzeigen. 1900, p.~742\,f.} We experience will's motives and will's results, but
not will itself --- just as in the domain of material nature we experience
energy's conditions and expressions, but not energy itself. \textsc{Hume} has
already shown this for all causality, and so for psychical and material
causality alike. The concept of will is formed, like the concept of energy, by a
construction. But it is a construction we are compelled to make. And if one
determines the concept of a psychical element, not in such a way that it must
necessarily be something that can be an object of immediate self-observation,
but in such a way that it designates an essential and irreducible property of
conscious life, then will may very well be a psychical element, and the concept
of will a fundamental psychological concept. The reason we cannot make will the
object of self-observation, as sensations, representations and feelings can be,
might lie in its extending through all the changing states and forms of
conscious life as their constant presupposition. Consciousness subsists only
through an unbroken labour of gathering the single elements into a whole. Such a
labour of combination and concentration expresses itself in the simplest
sensation as well as in every process of representation, every feeling, every
impulse and every decision. At every point there expresses itself an activity
which is just as original a side of conscious life as are the elements (sides or
properties) that observation and analysis immediately find before them. It is
not the case that we first have sensations, representations and feeling, and
that then, by their combination, something arises which we can call will.
Without an original combination, a primary synthetic process, the very elements
which determine will in the narrower sense do not arise at all. In the special
expressions of will (reflex action, impulse, wish, intention, decision) the
primary capacity of concentration expresses itself in a peculiar way under the
influence of certain determinate elements of knowledge and of
feeling.\footnote{See more on this in my \textit{Psykologi} II,~5; IV,~7~e;
V~B,~5; VII~A,~1; B,~4--5.} --- There thus takes place a constant interaction
between the element of activity, the will-element proper, and the intellectual
and emotional elements. We meet here again the antinomy mentioned above (1),
which asserts itself in all conscious life, indeed in all life. But the main
thing here is that we can form a concept of energy already on purely
psychological ground, since psychical work is performed wherever a psychical
phenomenon appears, this always --- so far as we can trace it ---
presupposing a synthesis. The psychical work in which the synthesis consists is
the greater, the more the single elements are qualitatively different, and the
further apart they lie in time.

It might now be tempting to identify this psychical energy straightway with the
energy that operates in the nerve tissue. But since not all nerve processes are
connected with phenomena of consciousness, a distinction must be drawn between
conscious and unconscious nerve energy, and the problem poses itself anew.
Besides, we know nothing more precise about this so-called nerve energy, or
about its relation to other forms of energy, organic and inorganic. In taking
the riddle to be solved --- as \textsc{Ostwald}\footnote{\textsc{Ostwald} seeks
in his \textit{Naturphilosophie} (Leipzig 1902) to show that the concept of
energy is the fundamental natural-scientific concept, and since the phenomena of
consciousness too must (with Kant) be conceived as expressions of activity
(\textit{Bethätigungen}), it becomes the fundamental psychological concept as
well. But about the relation between the two kinds of energy he does not express
himself clearly. Now the processes of consciousness are themselves called
``energetic'' (p.~394), now consciousness is said to be conditioned by the
nervous system's energy (p.~396), now spiritual energy is defined as ``conscious
\emph{or} unconscious nerve energy'' (p.~398).} and others do --- by the very
setting up of the concept ``nerve energy'', one merely introduces a
\textit{qualitas occulta} and acquiesces in it. The natural-scientific concept of
energy is, wherever one meets it, always abstracted from phenomena with
geometrical properties. Natural science knows energy only as an expression of
the relation between phenomena in space. The riddle would therefore be solved
only if we could form a concept of energy from which the psychological and the
natural-scientific concepts of energy could be derived as special forms; but we
lack the means to construct such a concept. In any case every attempt in that
direction would carry us beyond the domain of the psychological problem.

% ============================================================
%  II.  ERKENDELSESPROBLEMET  (pp. 28--53; notes 26--51)
% ============================================================
\fpchapter{II.\quad The Problem of Knowledge}
\label{ch:erkendelse}

\noindent\textbf{1.}~Not only psychological understanding, but all
understanding, is conditioned by the relation between continuity and
discontinuity. The latter conditions the fullness and the manifoldness of
understanding's content; the former conditions its gathering-together and its
ordering. Understanding appears under different principal forms, to which
different kinds of science correspond.

I understand what something is when I recognize it again. Thus I understand, or
know, who is coming when I recognize the one who is coming. Recognition rests on
there being a connection between new and earlier experiences. In the act of
recognition all intervening experiences are pushed aside, and the new phenomenon
is identified --- immediately or mediately, involuntarily or after some
deliberation\footnote{On the different kinds of recognition see my
\textit{Psykologiske Undersøgelser}, p.~35\,f.} --- with a phenomenon that has
occurred earlier. The contrast with what is new co-operates here as well; the
quality of familiarity becomes the more strongly marked, the more unfamiliar the
surroundings are. --- It is recognition that makes \emph{description} and
\emph{classification} possible; though the apprehension of difference
co-operates essentially in both operations. In all describing and classifying
certain types stand fast, at least provisionally fast, and the new phenomena are
referred to them, either as entirely the same (coinciding with one another), or
as similar (qualitatively alike), or as presenting analogy (likeness of
relation). The positive concepts we form express such types. Where we can have
recognition of no degree at all, and so can set up no type, we provisionally
collect the phenomena under a negative concept --- that is, we give them the
common mark of being different from the types hitherto determined. The negative
concepts have this significance in the history of classification, that by their
help one gathers what has hitherto been undetermined into a group distinct from
what is determined, so as then to leave it to later research to find positive
determinations within it.\footnote{See \textit{Det psykologiske Grundlag for
logiske Domme}, p.~32\,f.}

In the history of philosophy Plato's doctrine of Ideas stands as the
characteristic expression of the significance of describing and concept-forming
inquiry. The very possibility of recognition amid the motley manifoldness of the
phenomena fills Plato with enthusiasm, and the highest spiritual capacity was
for him that by which the different types underlying recognition are ordered
according to their mutual likeness and difference, so that a world of thought is
formed within which one can ascend and descend by steps that form a continuous
series.

When concepts have been formed in different connections, they can be combined in
judgments, and when different judgments have concepts in common that can be
substituted for one another, inferences can be drawn. By the road of
\emph{inference} concepts and judgments can be formed without one's needing to
go back to experience, as the describing and classifying sciences are obliged to
do at every point. A continuity of a higher kind here asserts itself, in that
thoughts and series of thoughts of the most various origin can be brought into
connection with one another. While recognition rests merely on identity,
inference rests on rationality: it is the relation between ground and consequent
that governs, and that makes a new kind of understanding possible. I understand
\textit{A} = \textit{C} when I know \textit{A} = \textit{B} and \textit{B} =
\textit{C}. This kind of understanding is won through the formal sciences. It is
peculiar to them that in the end it is a matter of indifference whence the first
concepts and judgments derive, provided only they are of such a character that
chains of inference can be built upon them. Just so, the describing sciences do
not trouble themselves about where the phenomena come from, provided only they
can be ordered according to relations of likeness and difference.

A third kind of understanding is characterized by this, that one neither merely
ascends from phenomena to concepts, nor merely forms by inference new
combinations of concepts out of given combinations of concepts, but derives new
phenomena from phenomena given earlier. The new phenomena are understood by
being set in a relation to earlier phenomena analogous to that which subsists
between ground and consequent in an inference. This kind of understanding is
characteristic of the exact real science (natural science), as it has developed
especially since the age of the Renaissance. It is here the \emph{concept of
cause} that governs. We here connect experiences according to their unavoidable
and law-determined succession in time.

This third kind of understanding has especially occupied the epistemology of
more recent times, after \textsc{Hume} and \textsc{Kant} had impressed the
difference between logical inference and real causal explanation. The
possibility of forming and applying the concept of cause has in more recent
times awakened a wonder similar to that which the possibility of forming general
concepts awakened in Plato. It is, however, not only the intellectual need to
find connection among experiences that has led to placing the concept of cause
first, but also the need to distinguish sharply between subjective
representations and objective reality, since in doubtful cases the criterion of
reality is always in the end the firm, unbreakable connection.\footnote{See
\textit{Det psykol.\ Grundlag}, Kap.~6.} The world of dreams and of thoughts is
larger than the world of reality (Eng ist die Welt, und das Gehirn ist weit),
and men have felt ever more the need to get their thoughts so determined and so
connected that they can stand as the expression of a reality --- especially
since only thereby can they hope to find the means of attaining their ends.

The concept of cause appears in two forms: a provisional, elementary form at
which we are often compelled to stop, and an ideal form towards which all
inquiry and all theories strive. \emph{The elementary concept of cause}
indicates merely an unavoidable relation of succession: when the phenomenon
\textit{A} has occurred, the phenomenon \textit{B} thereupon unavoidably occurs,
and \textit{B} occurs only when \textit{A} goes before. It is not thereby said
that the causal relation obtains between \textit{A} and \textit{B} themselves;
it is possible that both are successively appearing effects of an antecedent
cause. \emph{The ideal concept of cause} goes a step further and sees in the
phenomenon called the effect a continuation of, or a conversion into new form
of, the phenomenon called the cause. The ideal concept of cause thus passes over
into the concept of development, wherefore it is no wonder that this concept,
alongside the concept of cause, plays so great a part in the science of more
recent times.

\noindent\textbf{2.}~All three kinds of understanding rest on certain
fundamental propositions or principles which involuntary thought-work --- as it
is practised by practical common sense, in the special sciences, in
philosophical speculation and in religious thinking --- feels no need to seek
out, but which epistemology brings forward and tests. The problem of knowledge
arises when it is asked what \emph{the validity of understanding} rests on, and
how far it extends. It is the relation between continuity and discontinuity that
here again poses the problem. For Plato the problem arose from the irrational
relation between concept and phenomenon; in more recent times, from the
irrational relation between formal and real science. Can the concept ever be an
adequate expression for the manifoldness of the phenomena, and the law for their
changing combinations?

For \textsc{Plato} the consequence was that the concepts (``the Ideas'') must
have been known to man from a higher existence, and now emerged for thought in
their whole perfection under the influence of the imperfect imitations of them
which the world of sense presents. In more recent times the thought of nature's
lawfulness has often been taken for an a priori truth, an original intuition
which at most owed to experience the occasion of its being called forth. It is
the incommensurability between the principles (the logical principles and the
principle of causality) and experience that has given birth to this speculative
epistemology. This incommensurability has also given birth to a theory one might
call the arbitrary theory, because it strongly stresses what is arbitrary in the
first setting up of the principles --- but which stresses no less strongly that
the principles in their pure form can win no confirmation from experience, can
never become results, but always only postulates. There can --- says
\textsc{Hobbes}, the most pronounced representative of this standpoint --- be no
science of the very principles of all science; they are set up by constructive
art (\textit{principia sunt artis sive constructionis, non autem scientiæ et
demonstrationis}), and it is therefore we ourselves who bring forth their truth
(\textit{rationis prima principia vera esse facimus nosmet ipsi}). In
\textsc{Fichte} the first and unconditioned fundamental proposition of all human
knowledge is set up by a free act; our science rests not on a ``Thatsache'' but
on a ``Thathandlung'', namely the thinking consciousness's resolve always to
remain in agreement with itself: the setting up of the principle of identity.
Similarly \textsc{S.~Kierkegaard} assumes a leap in the setting up of the
principles, and \textsc{Kroman} derives the principle of identity from the
instinct of self-preservation, which seeks to maintain consciousness's unity;
``from experience it has in no way arisen''.\footnote{\textsc{Hobbes}:
\textit{Logica.} Cap.~3 (§§~8--9); \textit{Physica.} Cap.~25 (§~1). (Elsewhere,
however, Hobbes teaches that it is analysis proceeding from what is given that
leads us to the principles: \textit{Analytica est ars ratiocinandi a supposito
ad principia, id est ad propositiones primas vel ex primis demonstratas, qvot
sufficiunt ad suppositi veritatem vel falsitatem demonstrandam. De rationibus
motuum et magnitudinum.} Cap.~20 (§~6).) --- \textsc{Fichte}: \textit{Grundlage
der gesammten Wissenschaftslehre.} §~1. --- \textsc{S.~Kierkegaard}:
\textit{Uvidenskabelig Efterskrift}, p.~83. --- \textsc{Kroman}: \textit{Vor
Naturerkendelse}, p.~270--274.}

Both the speculative and the arbitrary theory must, however, concede that there
must be determinate empirical occasions which call forth the intuitions or the
postulates. Empiricism (which --- when it takes account not merely of the single
individual's experiences but of the whole race's --- passes over into
evolutionism) now lays, where it unfolds in its absolute form, the chief weight
on these ``occasions'' and regards them as the complete causes. It is most
clearly presented by \textsc{Stuart Mill} and \textsc{Herbert Spencer}. It seeks
to show how general principles may have formed themselves successively in
consciousness under the constant influence of the surroundings, and it
energetically maintains that no principle can have any validity beyond the
confirmation by experience (verification) that it can win. Only as results, then,
have the principles value.

What empiricism cannot explain is the fact that the principles themselves call
forth experiences, through the questions they lead us to ask. More happily
placed here is a fourth theory, developed in the most recent period by
\textsc{Ernst Mach} and \textsc{Richard Avenarius}, which one might call the
economic. It lets both passive experience and active development of thought come
into their right, in that it regards the principles as those ``conceptual
reactions'' which by the shortest road make intuitive apprehension and survey
possible. When the principle of continuity plays so great a part, it is because
it is a decidedly economic principle. There is here, then, only a question of
practical applicability. One forms the concepts and principles one needs in
order to make oneself master of the phenomena. Every practical and intellectual
need is satisfied when our thoughts are fully able to reproduce the sensible
facts. This reproduction is physics's end and aim, and atoms, forces and laws
are only means of easing the reproduction. Their value reaches only so far as
they help us.\footnote{\textsc{Ernst Mach}: \textit{Beiträge zur Analyse der
Empfindungen.} p.~144.} \textsc{Maxwell} and \textsc{Hertz} have expressed
themselves in a similar direction.\footnote{\textsc{Maxwell}: Scientific Papers.
II, p.~360. --- \textsc{H.~Hertz}: \textit{Einleitung.} (Werke III, p.~1\,f.).
--- Compare an interesting discussion \textit{De la valeur des lois physiques}
in Bulletin de la société française de philosophie. Mai 1901.}

The four theories, which respectively conceive the principles of knowledge as
intuitions, postulates, generalizations and economic instruments of thought, all
presuppose \emph{the analytic or regressive epistemology} which \textsc{Kant} in
particular has developed. For \emph{what} is to be the object of intuition, or
of postulate, or of generalization, and \emph{what} answers to the economic
needs of research, can be discovered only by \emph{inferring back from the data
which experience presents, and finding the presuppositions on which the
understanding of them rests}. Such an analysis must lie at the ground of every
epistemology. Whether this analysis has been completed, one will not be able to
convince oneself with full certainty. The history of science shows that here is
a labour that must be taken up again and again. Now one supposes one needs more,
now fewer principles than hitherto. A guarantee that one has reached the
absolutely last presuppositions one will never be able to obtain. What the
economic theory --- to which I shall chiefly hold --- particularly stresses is,
partly that no more principles are to be set up than what is given demands with
strict necessity (this is what is meant by ``critique of pure experience''),
partly that at different times, in different scientific situations, different
principles may be needed, so that even a principle which has for a long time
carried research forward may later be recognized as invalid, without its
historical significance being thereby misjudged. The economic theory thus
stresses partly the law of parsimony, partly the purposive, utility-determined
character of the principles. It is due especially to two classes of motive:
first, the wish to reduce to the least possible those presuppositions that
cannot themselves be proved --- an antidogmatic or antimetaphysical motive; and
secondly, experiences from the history of the sciences, which show how
principles and hypotheses may be legitimate and fruitful in a certain period but
must later be replaced by others --- something to which the most recent period's
discussion of the grounding and the validity of the mechanical view of nature
has particularly drawn attention.

There are, however, two sides of the matter which the special form the analytic
epistemology has assumed in appearing as economic theory has not sufficiently
stressed, and which become of particular importance when one does not regard the
problem of knowledge in isolation but sees it in its connection with the other
philosophical problems.

In the first place, the forms in which the intellectual need is satisfied must
hang together with the general nature of consciousness. \emph{What} we
understand, and \emph{that} we understand anything at all, depends not merely on
the constitution of the phenomena but on our intellectual organization, just as
which colours we see depends on the constitution of our organ of sight and not
merely on the outward objects. There will be a certain type for all principles
and hypotheses, one which points back in the end to the innermost nature of
conscious life. And it will here always be the need for unity and continuity
that one comes back to. Which and how many fundamental concepts (categories) and
fundamental presuppositions one must set up is a problem that must always be
taken up anew in knowledge's struggle to press forward. It was an illusion when
Kant supposed that one could give, once for all, a catalogue of what is needed in
this respect; but the old master was not mistaken in this, that it is \emph{the
need for unity and continuity} that lies at the ground of all the forms in which
understanding is won, or is supposed to have been won. He has himself shown that
all his categories can be traced back to the concept of continuity, and through
all vicissitudes in the domain of principles it will undoubtedly be this concept
that constantly recurs. The logical principles, the principle of causality and
the fundamental propositions of natural science all lead back to this concept,
which hangs so closely together with the nature of conscious life. A purely
psychological epistemology will, however, never be able to suffice; for from the
fact that the need for unity and continuity --- be it ever so essential to
conscious life --- is satisfied by certain principles, it by no means follows
that these are objectively valid. That need can find vent and satisfaction in
many ways and in many forms, as the history of mythologies and of speculations
sufficiently establishes --- ways and forms which as a rule by no means satisfy
the demands of economy, whether in respect of parsimony or in respect of
applicability. Here \textsc{Schiller}'s words hold: Eng ist die Welt, und das
Gehirn ist weit. Of all the wealth of thought at the human spirit's disposal,
only a narrow and firmly closed series of thoughts can be used when it is a
question of valid understanding. What is a necessary presupposition is only
this, that the presuppositions which the understanding of what is given requires
must be psychologically possible, must accord with the general laws of conscious
life, and be only more special and more exact developments of what lies in those
laws. A selection must take place among the thoughts that involuntarily press
forward; but it brings with it no emancipation from the general psychological
laws. What matters is that new dispositions of a special kind should form
themselves --- that there should arise an intellectual habitus which poses the
questions and criticizes the answers in a stricter and more determinate way than
the involuntary course of thought feels any need to do. Every comprehensive
principle is really --- psychologically regarded --- the expression of such a
habitus, which may be more or less deeply grounded in the nature of
consciousness, that is, may now have the character of an instinct, now be due to
the influence of practice. The purely logical principles come closest to the
instinctive. The need for agreement with oneself, for consistency in the course
of representations, cannot be explained from parsimony and expediency alone.
Where strict inference from the experiences so far available leads us to
self-contradictory results, we far sooner assume that the experiences are
incomplete than that existence should contradict itself. What holds in the
highest degree of the purely logical principles holds of more special principles
as well. Thus the principle of causality expresses a tendency, when an event has
occurred, to look about for other events in which the conditions of its
occurrence might be found. Here too a need asserts itself to bring unity and
connection into the content of consciousness. The more special and determinate
the satisfaction of this need is to be, the more the economic principle operates
in its two forms: as parsimony, which takes the straight road to the goal, and
as expediency, which follows roads that really lead to understanding. Examples
of this are to be had in Kepler's and Newton's demand for \textit{veræ causæ},
in the principle of inertia, in the principle of energy, and so on. What purely
empirically stands as a hypothesis becomes, epistemologically regarded, a
principle, a guiding thought under whose leadership consciousness can, in the
world of experience, satisfy its need for unity and continuity.

In the second place, principles and fundamental hypotheses --- even if in the
end they all stand as \emph{working} hypotheses in the service of intellectual
economy --- must not be conceived as wholly accidental or arbitrary in relation
to the existence which we mean to understand by their help. The very concept of
a working hypothesis points in two directions: partly --- as already shown ---
back to the nature of the thinking consciousness, since our consciousness cannot
carry out a function, be it ever so economical, that is wholly foreign to its
own essence; and partly outward, towards the existence to which the phenomena
that are to be understood belong. A working tool must fit both the hand that is
to use it and the object that is to be worked upon. Existence must then in
itself present sides that answer to the unavoidable presuppositions of our
knowledge, however much these are conditioned by the inner and outer
circumstances under which knowledge works --- or perhaps precisely for that
reason. This holds for all valid knowledge, from the most elementary forms to
the highest. It holds already of the sense-qualities, despite their
``subjectivity'', and it holds of the highest principles of unity in our
thinking. \textsc{Kant} says of the tendency to presuppose a unity behind the
differences of natural phenomena:\footnote{\textit{Kritik der reinen Vernunft.}
p.~653.} ``One might perhaps believe that we had here a merely economical
contrivance of reason, serving to spare effort~.... But such an egoistic purpose
is very easily distinguished from the Idea according to which everyone
presupposes that this unity of reason accords with nature itself, and that
reason here does not beg but commands, although it cannot determine the limits
of this unity.'' I add, however: even where reason commands (or asks, expects,
anticipates, postulates), it is obliged --- as all commanders are --- to adjust
its commands to the capacity for obedience that may be presumed to be present.
---

This last consideration leads in particular to making possible the connection
between the problem of knowledge on the one side and the problem of existence on
the other. It is all the more important to bring this out, since the analytic
epistemology --- like the economic, and in its way the arbitrary as well ---
leads to putting a new concept of truth in place of the current, naive one. The
significance of the principles is to guide our work of winning understanding.
Their \emph{truth} consists in their \emph{validity}, and their validity in
their \emph{working value}. That a principle is true means that one can work
with it; and that means, when the principles of knowledge are in question, that
by its help one can get forward in understanding --- in firm ordering and
connection of the phenomena. The concept of truth is a \emph{dynamic} concept,
in that it expresses a determinate way in which the energy of thought is
employed, and a \emph{symbolic} concept, in that it does not mean coincidence
with, or qualitative likeness to, an absolute object, but likeness of relation
(analogy) between the events in existence and human thoughts. The old naive
concept of truth, according to which a piece of knowledge was true when it
directly pictured or mirrored ``reality'', is untenable, and has been so ever
since the moment when men began to assert \emph{the subjectivity of the
sense-qualities}. That the sense-qualities were subjective was not, after all,
to mean that they were invalid, unusable for orienting us in the world. They
still stand as signs, signals, symbols, whose mutual succession we can interpret
as the expression of an objective series of events, although they cannot be
proved to be pictures of it. It stands similarly with the logical principles and
the other fundamental presuppositions of our knowledge. \emph{Critical
philosophy} led to the result --- a simple consequence of the very analytic
method it employed --- that the truth of the fundamental propositions can mean
only that they make scientific experience possible; they were themselves, after
all, found by analysis as the necessary presuppositions of such experience. A
comparison between our thoughts and an absolute world of things is not possible;
we can compare only thoughts and experiences with one another. This consequence
Kant himself did not see so clearly as some of his disciples did (Maimon and
Fries);\footnote{\textit{Den nyere Filosofis Historie} II, p.~114\,f.; 222.} the
master himself was still entangled in dogmatism, as is seen from his doctrine of
the ``Ding an sich'' as absolutely subsisting outside every subject. But when he
designates the law of causality an ``analogy of experience'', meaning thereby
that events in time stand to one another as ground and consequent do in our
thinking, he comes close to making it a working hypothesis. In the most recent
period one is --- as we shall soon see --- \emph{on the side of natural science}
more inclined to acknowledge the dynamic and symbolic concept of truth than one
was so long as the mechanical view of nature stood with a certain dogmatic
character. --- This new concept of truth, which is pressing forward in the
domain of science, presents a kinship with the way in which --- as we shall see
under the fourth principal problem --- the religious consciousness more and more
places itself over against the dogmas; in the religious domain too it is more
and more applicability and working value that are asked after. The static
concept of truth\footnote{The expression ``static concept of truth'' was used
--- with a somewhat different motivation --- by \textsc{Louis Weber} at the
philosophical congress in Paris in 1900.} is everywhere being pushed aside by
the dynamic.

Even supposing that fruitful principles or working hypotheses have been won,
will existence then be exhaustively reproducible by their help, or will there
always continue to subsist \emph{an irrational relation between the principles
our consciousness is able to form and existence itself, from which our
experiences derive}? --- It will appear that under three forms especially
something irrational will always remain: in the relation between quality and
quantity, in the significance the time-relation has for the concept of cause,
and in the relation between subject and object. We shall consider these three
points each by itself.

\noindent\textbf{3.}~In the attempts to reduce all given differences to identity
and continuity, the endeavour to trace differences of kind back to differences
of degree has been of particular importance. In the science of material nature
it expresses itself as an endeavour to explain all change as motion in space.
Motion from one place to another is the simplest change; it would therefore be a
great advance in clarity if it could be shown to be the change which, under
different forms, takes place everywhere in nature where immediate experience
shows us qualitative changes. \textsc{Aristotle} already taught that spatial
motion lay at the ground of all other change, and of all coming-to-be and
passing-away as well; he did not, however, hold it possible --- as the Atomists
had wished --- to trace everything in material nature back to it. Only with the
rise of modern natural science does this thought come forward in earnest.
\textsc{Galilei} stands as the great founder of what one is accustomed (in the
narrowest sense) to call \emph{the mechanical view of nature}. Besides the
simplicity thereby attained --- and Galilei was certain that nature follows the
simplest ways --- there is gained the further advantage that one can operate
with determinate quantities. Galilei therefore said: measure all that is
measurable, and make measurable what is not! By this reduction agreement becomes
possible, whereas qualitative differences can only be estimated and not
measured; and exact verification becomes possible.\footnote{\textit{Den nyere
Filosofis Historie} I, p.~160--167.}

The principles on which the mechanical view of nature rests were regarded by its
most important representatives as absolute truths, fundamental laws of
existence; in so far as it was still thought that they needed a grounding, this
was taken directly from God's essence and will. In this Cartesians and
Newtonians were agreed. And the Materialists differed from them only in holding
the theological grounding superfluous and impossible. Through a series of
magnificent discoveries and explanations this view of nature has shown its
fruitfulness. The question for us is whether, from an epistemological point of
view, it can be said to mark an absolute conclusion.

Even supposing that everything in material nature could be explained by the
principles of the mechanical view of nature --- and there is, as we shall soon
see, doubt of that in the most recent period within the circle of natural
scientists --- it is in the first place clear that the qualities are not got out
of the world by being ``reduced'' to quantities or charged to the account of the
sensing subject. They remain standing as immediate facts which must be
established empirically. The chemical product's properties cannot be derived
from the properties of the elements; and when one kind of physical energy
receives its equivalent in another kind of physical energy, the equivalent
nevertheless has new properties which cannot be derived from the properties of
the first form of energy. One cannot, however, let the sensing subject create
these qualities out of nothing; in any case we then get insoluble psychological
difficulties in place of the physical and chemical difficulties we have pushed
aside.

In the second place, extension and motion are in the end themselves qualities in
the widest sense --- factual properties which in and for themselves might just
as well need an explanation as the specially so-called sense-qualities. Since
\textsc{Berkeley}'s and \textsc{Leibniz}'s time this has often been urged from
the philosophical side. There is no reason to assume that the quantitative
properties express the innermost essence of things. The reason science seeks
them out by preference and holds to them is really only that by their help one
can give exact descriptions of the material phenomena and draw determinate
inferences from them. This has in the most recent period been strongly
emphasized by investigators such as \textsc{Maxwell} and \textsc{Hertz} --- in
the latter's case with express adherence to the economic epistemology discussed
above. ``The progress of the exact sciences'', says Maxwell, ``depends on the
discovery and development of appropriate and exact ideas, by whose help we can
form a mental representation of the facts which is at once sufficiently
comprehensive to stand for every particular case, and sufficiently exact to
afford security for the inferences we draw from them by way of mathematical
calculation.'' And Hertz says: ``We form for ourselves such inner images
(\textit{Scheinbilder oder Symbole}) of outward objects that the logically
necessary consequences of the images are always in turn images of the naturally
necessary consequences of the objects imaged~.... Any other agreement with the
things, the purpose of our representations does not
require.''\footnote{See Note 31.} --- The dynamic and symbolic concept of truth
is here expressly put in the place of the naively dogmatic one which the
mechanical view of nature had formerly espoused. What matters is only to find a
circle of symbols that can be applied with full consistency, and from which
inferences can be drawn that are confirmed by new experiences, which in their
turn can themselves be expressed by the same circle of symbols. But the
possibility can here never be ruled out that another circle of symbols would be
able to express the positive experiences equally well or better, and to serve as
the ground of exact and verifiable deductions. It will never be possible to
prove of any circle of symbols that it is the only right one, and the only
necessary or possible one.

The epistemological reflections to which recent natural scientists have thus
been led were called forth by the difficulty that presented itself in bringing
the electrical phenomena under the mechanical view of nature, or else in
deriving that view's principles from the laws holding for those phenomena.
\textsc{Hertz} inclines rather to the latter possibility, \textsc{Boltzmann}
rather to the former.\footnote{\textsc{H.~Hertz}: \textit{Ueber die Beziehungen
zwischen Licht und Elektrizität.} 1889, p.~29. --- \textsc{Boltzmann}'s lecture
at the German congress of natural scientists 1900 (translated in The Monist.
January 1901: \textit{The recent development of method in theoretical Physics})
and his paper in the Vienna Academy's Oversigter CV, 8 (translated in The
Monist. Octbr.\ 1901: \textit{On the necessity of atomic theories in Physics}).
--- On the whole question compare also \textsc{C.~Christiansen}: \textit{Den
elektromagnetiske Lysteori.} (Det danske Vid.~Selsk.s Oversigter. 1889), and the
same author in \textit{Fysisk Tidsskrift} I, p.~4--5.} But of greatest interest
for epistemology is \textsc{Maxwell}'s criticism of the current concept of
matter, according to which by matter one imagines an extended mass. The fault in
this concept is, according to Maxwell, that it leads one to think of matter as
something at rest. But it is motion that makes rest intelligible to us, not the
reverse. The doctrine of motion must therefore precede the doctrine of
equilibrium, dynamics precede statics. It is from motion that we win the concept
of force or energy, by whose help equilibrium becomes intelligible to us. If now
by matter we understand what is constant, what subsists through all changes,
then it can only be the law of motion itself, and the essence of ``matter''
consists accordingly in motion. It is, according to Maxwell, a prejudice to
regard matter as extended and the molecules as hard and solid: one would then
have to ask what holds the molecule's parts together, and one would arrive at
molecules of a second order. Yet we must constantly operate both with
geometrical and with dynamic concepts, whether we regard the last elements as
extended or not. It is only extension as a resting property that Maxwell
declares against; even with a line we draw on a board, the motion by which it
comes to be is the essential thing.\footnote{\textsc{Maxwell}: \textit{Scientific
Papers.} II, p.~26--33; 302--305 (where M. shows what he owes to Faraday and
Lord Kelvin); 326; 777\,f.} --- It is an extraordinarily important
epistemological point of view that is here brought forward: that resting states
always contain problems which can be solved only by rest's being replaced by
motion. Potential energy we understand only through actual energy, capacities
and dispositions only through what comes of them. This law holds everywhere, in
the psychical domain as in the physical. There is a half mystical, half
materialistic inclination to find completion in contemplation of --- or staring
at --- a resting ``something''.\footnote{Compare my \textit{Psykologiske
Undersøgelser}, p.~77--81.} Maxwell has made an important contribution towards
getting this inclination rooted out. But the great question is whether the
thought of the continuity of motion, or of activity, can be carried through in
all domains. Even if the dynamic asserts, epistemologically regarded, its
priority over the static, the static cannot for all that be got rid of; it
presents itself ever anew with its tasks for thought, and Maxwell himself
acknowledges it when he declares both geometrical and dynamic concepts necessary
in the explanation of nature. In contrast to the dynamic, the geometrical
designates simultaneity. In the domain of material nature the simultaneous
appears in the form of space;\footnote{This point \textsc{Ostwald} does not let
come into its own when he wants to trace everything in nature back to
``energy''. Mass, the filling of space, weight and chemical properties are all
supposed to be traceable back to the concept of energy. It is correct that all
those concepts presuppose the concept of energy; but is that the only
fundamental concept in natural philosophy? In that case one would have to be
able to derive the ``geometrical'' properties from the dynamic ones, and at that
Ostwald makes no attempt. In his work \textit{Die Ueberwindung des
wissenschaftlichen Materialismus} (1895) he even defines matter as ``a spatially
[sic] ordered group of different energies'' (p.~28). This spatial ordering, here
acknowledged as a special moment, is --- it seems to me --- somewhat pushed
aside in his \textit{Naturphilosophie} (1902), where ``the energetic
world-picture'' is developed. Yet his main proposition, on which he builds his
system, is that ``every process without exception can be exactly and
exhaustively described or presented by indicating which energies undergo
temporal or spatial [sic] changes'' (p.~152). That the concepts energy and time
cannot be separated is a matter of course, since energy means the capacity to
overcome a resistance, and every overcoming takes time. Not so much a matter of
course is the connection between the concept of energy and spatial changes. This
connection derives, in Ostwald, simply from his having abstracted the concept of
energy from spatial phenomena, from investigations of changes at determinate
places in space. Only by this geometrical moment's being pushed aside in his
concept of energy does it become possible for him to believe that the concept is
without further ado the same for psychical and for physical phenomena, so that
the problem of ``soul and body'' would be solved by setting it up. See above,
Note 25.} in the psychical domain the relation of simultaneity does not assume
this geometrical form, but appears nonetheless as something ``static'' which
sets research new tasks when it has happily and successfully found the laws of
the psychical processes. After the continuity in the processes of development
has been exhibited, it remains to exhibit that there is continuity between the
processes and the resting states. --- We meet here again, then, the standing
problem; it would present itself even if the qualities were fully explained by
being represented by quantities, and even if the fundamental physical concepts
with which science operates were more than the most perfect and most expedient
circle of symbols that it has so far succeeded in forming and applying. The
possibility of an irrational relation between existence and our knowledge cannot
therefore be ruled out.

\noindent\textbf{4.}~The investigation of the relation between \emph{the
elementary and the ideal concept of cause} will lead us to a similar result. In
its elementary form the concept of cause indicates, as we have seen, only this
relation between two events: that when the one has occurred, the other
unavoidably occurs. But from this unavoidable succession research advances,
through observations, experiments and hypotheses, so far as possible to the
exhibition of a continuity in the series of events, by which the differences
between the single members of that series become vanishing, so that at last even
the temporal difference between the events is reduced to the least possible.
From outer, though unavoidable, succession thought works its way through
continuity towards a complete identity. While the time-relation plays an
essential role in the elementary concept of cause, it is as good as eliminated in
the ideal concept of cause. If this process could be carried through completely,
we should arrive at the paradoxical result that the perfect causal explanation
would be an abolition of the concept of cause itself: for the causal relation
differs from the purely logical relation between ground and consequent only by
the temporal difference between the members of the relation.

The philosophy of the seventeenth century (\textsc{Spinoza} and
\textsc{Leibniz}) did not distinguish between cause (causa) and ground (ratio);
the causal relation between two phenomena was for them the same as the relation
between premisses and conclusion in an inference, a mere relation of identity. It
stood for them as self-evident that there could not be more in the effect than
there was in the cause; the time-relation between the members was indifferent,
and time on the whole belonged only to the obscure sphere of sensuous
representation (imagination). When \textsc{Hume} raised the problem of causality,
he did so precisely by emphasizing that no proof can be given for the
justification of inferring from past to future. When appeal is made to the
experiences of the past, Hume says, these can only show that a thing has at some
time, in a single moment, been in possession of a certain power, not that it
constantly possesses it and will possess it.\footnote{\textsc{Hume}:
\textit{Treatise on human nature.} I, 3, 6: Your appeal to past experience
decides nothing in the present case, and at the utmost can only prove, that that
very object, which produc'd any other, was \textit{at that very instant} endow'd
with such a power.} \textsc{Kant} believed indeed that he could give a rational
proof of the validity of the causal principle, but he distinguished between the
causal relation and the relation between ground and consequent, in such a way
that there was only an analogy, not an identity, between them: the law of
causality meant for him that events follow one another in time in a manner
analogous to the manner in which the conclusion springs from the premisses. The
causal principle therefore gives us a rule by whose help we can attain unity in
our experiences.\footnote{\textit{Kritik der reinen Vernunft}: p.~181: Wir werden
also durch diese Grundsätze [the causal principle and the kindred principles
concerning the persistence of ``substance'' and the reciprocal action between all
that exists] die Erscheinungen nur \textit{nach einer Analogie mit der logischen
und allgemeinen Einheit der Begriffe} zusammen zu setzen berechtigt werden. ---
p.~180: Eine Analogie der Erfahrung wird also nur eine Regel sein, nach welcher
aus Wahrnehmungen Einheit der Erfahrung entspringen soll, und als Grundsatz von
den Gegenständen (der Erscheinungen) nicht constitutiv, sondern nur regulativ
gelten.}

From two different standpoints attempts have been made in the most recent time
to eliminate the time-relation and thereby to abolish not only the elementary
concept of cause, but in the end the concept of cause altogether.

From a speculative standpoint it has been maintained that the time-relation
always points to an imperfection, to an incompleteness in development. So long as
the time-relation is determinative for our conception of existence, our
conception is therefore always imperfect, and the reality of the time-relation is
irreconcilable with the idea of a perfect knowledge, of an absolute truth. Our
knowledge does in fact work itself more and more away from the time-relation; the
more clearly we understand something, the smaller the role the temporal
difference plays --- the more does real knowledge pass over into formal
knowledge, cause into ground. What we begin by calling cause and effect comes, as
knowledge advances, to stand as members in a whole, members which stand in a
lawlike relation to one another, a rational relation for which the temporal
sequence is unessential. That we first observe the one member and thereafter the
other has no significance: in our knowledge of the law the whole process stands
before us as a unity. It is not the time-relation, but the unity behind the
time-relation, that binds what we call cause to what we call effect. Besides,
neither is there any interval between the end of the event we call cause and the
beginning of the event we call effect. It is the English philosophers
\textsc{Francis Bradley} and \textsc{Bernard
Bosanquet}\footnote{\textsc{Francis Bradley}: \textit{Appearance and Reality.}
London 1893. Chap. 4--7; 18. --- \textsc{Bernard Bosanquet}: \textit{Logic or
the morphology of knowledge.} Oxford 1888. Book I. Chap. 6.} who have urged this
consideration, at the base of which lies a speculative interest in a total
conception and in concluding concepts.

From a wholly different standpoint, which may be designated the standpoint of
``pure experience'', and whose most important representatives are
\textsc{Richard Avenarius} and \textsc{Ernst Mach}, work is done in a similar
direction, but from other motives. The interest here is directed to separating
out all the accessory representations and auxiliary concepts with which we,
involuntarily or on methodological grounds, supplement experience itself, the
immediately given. The advance of knowledge consists in a reduction of the
differences (to a ``heterotic minimum'') and in an approximation to a pure
description of a continuous process. In the course of this advance the content of
knowledge is more and more restricted to descriptive statements, so far as
possible with analytic transitions, and the differences are reduced from
qualitative to quantitative, so far as possible with a constant exhibition of
relations of equivalence. It is no longer the task of strict science to
``explain''; whoever wants ``explanations'' is referred to mythology and
metaphysics; the aim of science is to give an exact, methodical description of
all relations and transitions. As regards the causal relation in particular,
weight is laid partly on this, that every application of the concepts cause and
effect arbitrarily draws two elements out of the whole connection in which they
occur and sets them in opposition to one another; partly on this, that the
time-relation may very well be reversed as soon as one has exhibited a relation
of equivalence, and as soon as one abstracts from the direction of the
changes.\footnote{\textsc{Rich. Avenarius}: \textit{Kritik der reinen Erfahrung.}
II. p.~332--339. --- \textsc{Ernst Mach}: \textit{Zur Analyse der
Empfindungen.} p.~168: ``Die Thatsache der Nichtumkehrbarkeit der Zeit reducirt
sich darauf, dass die Werthänderungen der physikalischen Grössen in einem
\textit{bestimmten Sinne} stattfinden. Von den beiden analytischen Möglichkeiten
ist nur die eine wirklich. Ein metaphysisches Problem brauchen wir hierin nicht
zu sehen.'' --- Compare on these theories \textsc{Heinrich Grünbaum}:
\textit{Zur Kritik der modernen Causalanschauungen} (Archiv für systematische
Philosophie. 1899. p.~392--409).}

The epistemological views which thus, with different motivations, would
eliminate the elementary concept of cause and deny the fundamental significance
of the time-relation may be characterized by this, that they put an ideal of
knowledge in the place of the actual knowledge we can attain at any given time.
Every determinate investigation must begin at a determinate point, which lies
where the problem presents itself. And the problem is set when two members in the
series of events draw attention to themselves and awaken surmises of an inner
connection between them, or when a single member steps forth in its peculiarity
or its suddenness and thereby awakens the need to seek intermediate members
through which it can be brought into connection with the whole remaining series.
It is to that extent an accidental, arbitrary starting-point that is seized upon;
but from the nature of the case it cannot be otherwise. Our knowledge develops
historically, because human attention is awakened only under certain determinate
conditions. If it were directed without difference and at all times equally upon
all members in the series of events, no knowledge would arise. It is of course of
importance that one be conscious of what is accidental or subjective in the
starting-point and in the section taken; and if the work of knowledge has
success, one does get beyond it, since the section is precisely through
understanding worked into a great continuous connection.

The same historical character of our knowledge becomes clear when we bethink
ourselves that we always stand between the experiences of the past and the
possibilities of the future. \emph{In each single moment} we must distinguish
between what lies given and what we expect; the former is the cause, the latter
the effect; and the expectation can never become more than a hypothesis. The
single moment, in which on the one side a ``no longer'' holds and on the other
side a ``not yet'', sets the problem before us in its whole tension, a tension
which only the power of habit diminishes. With new problems the tension will not
fail to appear --- and so long as knowledge goes forward, it will seek and find
new problems.\footnote{\textsc{Bosanquet} accordingly concedes this too:
\textit{At any given moment} we have no choice but to say that the future is
conditioned by the past, \dots\ effect by cause. \textit{Logic} I. p.~272.} The
full connection is always present only afterwards, and until it is present the
assumption of it stands as a hypothesis. We shall therefore never be able to
dispense with concepts such as force, energy, cause or possibility, which with
different nuances express one and the same relation, namely the dependence of
later states on the preceding states; this dependence differs from purely logical
and mathematical relations of dependence in that it is at once a time-relation
and a rational relation, since the following state both follows after and follows
from the preceding one. Even if the continuity be exhibited by the help of ever
so many intermediate members and nuances, this relation continues to hold for
every little transition between two members or nuances.

And by this nothing is altered, even if equivalence is exhibited between the
states that succeed one another. Hume's problem has by no means, as it has
sometimes been said, been solved by the discovery of the conservation of energy.
A relation of equivalence does not exclude a qualitative difference, but
precisely presupposes it: I have no ground for setting up the judgment
\textit{A = B} when \textit{A} and \textit{B} have not first stood before me as
different, before I found on closer investigation that they could be put in the
place of one another. When such a relation of identity has been found, it is
indifferent whether we go from \textit{A} to \textit{B} or from \textit{B} to
\textit{A}; but at the time when it was to be found, we began with \textit{A} or
\textit{B} and came by way of investigation over to the other concept. In a
judgment we must therefore distinguish between the psychological process by which
the judgment comes to be, and in which a determinate difference asserts itself
between the initial representation (subject) and the concluding representation
(predicate), and the finished judgment, which can be formulated as a relation of
identity in which the difference between subject and predicate becomes
indifferent.\footnote{See \textit{Det psykologiske Grundlag for logiske Domme.}
p.~46--50.} And just as little as in the movement of thought is the order of the
members indifferent in outer happening. If I have recognized that there is a
relation of equivalence between heat and motion, it is, purely abstractly
regarded, indifferent whether I go from heat to motion or from motion to heat;
but ``the direction of the changes'' is in reality a question of life and death:
existence actually becomes another thing according as the changes proceed
preponderantly in the one or in the other direction. Time then cannot in reality
be reversed, and therefore the exhibition of relations of equivalence cannot form
the conclusion of our knowledge. We must in addition have information about the
actual direction of the changes (the conversions), and that is won only by ever
new experience.

There comes in addition that every relation of equivalence --- as well as every
causal relation whatever --- comes into force only when certain conditions, and
particularly releasing conditions, are present. With respect to the presence of
these conditions a new problem then arises: can it be shown, with respect to them
too, that the relation to their cause may be presented as a relation of
equivalence? We come here into an infinite series, where the answers (the
determinate, concrete answers) cease long before the questions cease. Here too we
have a purely logical analogy; for every absolute relation of identity
(\textit{A = B}) between two concepts rests, on closer consideration, on
determinate conditions (so that it --- according to \textsc{Jevons}'s
formulation --- must be expressed by \textit{AC = BC});\footnote{\textsc{Jevons}:
\textit{Principles of Science.} 2. ed. London 1877. p.~43. (Already
\textsc{De Morgan}, as Jevons remarks, has said that we usually think and argue
within a limited world or sphere of concepts, even when this is not expressly
declared.) --- Compare my \textit{Formel Logik} 3, \S\,16 and \S\,22.} it then
becomes a new question how matters stand with these conditions (\textit{C}), and
so thought goes ever further. ---

There is thus no prospect of escaping the historical element in our knowledge.
The measure of the development of knowledge consists in the extent to which the
elementary concept of cause (unavoidable succession) can be applied, as compared
with the mere ascertaining of actual simultaneous or successive differences, and
next in the extent to which the elementary concept of cause can be replaced by
the ideal concept of cause (equivalence or identity). But the process of
knowledge will at every time consist in an ascent through the three stages here
indicated --- a process of ascent which must ever anew be repeated from new
starting-points. And it is the reality of the time-relation that conditions this
necessity. Speculative conceptions --- whether they appear in idealistic or in
realistic form --- have therefore always an inclination to push the time-relation
aside or to regard it as merely ``empirical'', if not as
illusory.\footnote{\textsc{Francis Bradley} has clearly seen the contradiction
between time and ``the Absolute'' when he declares: If time is not unreal, I
admit that our Absolute is a delusion (App. and Reality: p.~206). (The Absolute
has no seasons. ib. p.~500). --- The critique of, and the reaction against, a
speculative tendency therefore also often supports itself on the reality of time
as a main argument. Compare \textsc{C. H. Weisse}'s and \textsc{S. Kierkegaard}'s
relation to Hegel (\textit{Den nyere Filosofis Historie.} II. p.~243; 261\,f.) In
his acute work \textit{Tidsexistensens Apologi. Et stycke relationsteori}
(Upsala 1888) \textsc{Pontus Wikner} has attempted to shake the either-or that
here presents itself. He wants to show that the highest perfection can be
attained only by a successive unfolding of qualities which as simultaneous would
exclude one another. But the limitation that lies in succession itself is not
thereby abolished!} Should it be an illusion, then it is in any case an illusion
in the second power when one believes one can free oneself from this illusion.
For us existence never dissolves wholly into thought.

\noindent\textbf{5.}~From yet a third side the problem of knowledge shows itself
in its whole sharpness, at the same time as a continual development of knowledge
shows itself possible. In every act of knowledge a distinction can be drawn
between \emph{a subjective and an objective element}, between that which knows
and that which is known --- but both elements are given only in relation to one
another, though within this relation they can assert themselves in different
degrees.

What now is subjective, and what is objective, in our knowledge? --- Already
from the foregoing it will be seen that this question is answered in different
ways. In the domain of science a tendency asserts itself to charge all
differences of quality, everything that breaks continuity, and at last perhaps
even everything that breaks identity, to the subject's account. The
sense-qualities and the differences of space and time would then be merely
subjective. What is justified in this way of looking at the matter rests on this,
that the differences we posit in qualitative, extensive, intensive and protensive
(temporal) respects depend on our psychical dispositions. Our sensation is a
discriminating whose result depends on the organization and the previous history
of the sensing subject. The differences found hold, therefore, only in relation
to the standpoint of the knowing subject in those various respects. Add to this
that our attention is discontinuous, moves by leaps to and from its object; by
that too there come forward differences and interruptions which cannot be charged
to the object's account.

The dogmatic and speculative tendency in philosophy and in natural science has
been inclined to take this road. Emancipation from the merely subjective became
here at the same time an approximation to continuity and identity as the essence
and the norm of truth.

The critical philosophy emphasizes, in contrast to this, that the qualitative,
extensive, intensive and protensive differences make up the material given for
our knowledge and set us the tasks of our inquiry. Just as our personality works
at gathering its sporadically appearing elements, at harmonizing the conflicting
tendencies and freeing itself from what is obscure and self-contradictory, so our
understanding works at transforming the given differences, which for us are
matters of fact, into stages in one and the same continuous process of
development, or into forms of one and the same content. The need for continuity
and identity lies deep in human consciousness, and consciousness therefore seeks
to find them again in the given content of knowledge. Moreover, consciousness
cannot itself bring forth the differences which form the material of its working,
however much the form and the degree in which they appear in consciousness are
determined by consciousness's own involuntarily operating conditions. --- In the
discussions of the most recent time concerning the epistemological foundation of
natural science there shows itself, too, a tendency to derive all points of view
of unity preponderantly from the economy of the knowing subject.

In reality we never at any time have the pure subject with its forms given in
opposition to a pure object, or rather to a ``thing in itself'' from which the
manifoldness of the content of knowledge would derive. \textsc{Kant} concluded
here too quickly, in supposing that the subjective forms of knowledge could be
determined once for all, so that the ``matter'' in knowledge would have to be
derived from ``das Ding an sich''. In the more special development of his theory
of knowledge he could not, however, avoid coming to ask whence the forms
derived, and to hint that these too (which are, just as much as the matter,
factually given, since they must after all be found out by psychological
analysis) are ultimately determined by the ``Ding an sich''; and since the
``forms'' designate what is constant in our knowledge, the last presupposition of
Kant's philosophy becomes something for which he can really give no grounding,
namely that the ``Ding an sich'' works in a constant manner, since otherwise the
forms could not subsist, or in any case could not be applied.\footnote{See my
treatise \textit{Kontinuiteten i Kants filosofiske Udviklingsgang}, p.~26--35.}
Thus Kant comes to oscillate between the two views described above, though it was
plainly his intention to hold to the latter. Continuity was for him the common
fundamental form of the categories, and synthesis was for him the fundamental law
of the working of the knowing consciousness.

The problem is more involved than Kant saw. When we distinguish between subject
and object in our knowledge, it is in reality an objectively determined subject
(\textit{S}\textsubscript{\textit{o}}) that we set in opposition to a
subjectively determined object (\textit{O}\textsubscript{\textit{s}}). The
properties or ``forms'' we ascribe to the subject cannot be explained from the
concept of subject itself (the pure \textit{S}); they stand as facts, just as
much as all the other properties our knowledge gets to deal with. And the
properties or determinations we ascribe to the object belong to it always only in
relation to a subject, and, more precisely, to a subject of a certain special
constitution. Therefore the problem begins ever anew: whence does the subject get
its objective (factual) properties, and what relation is there between the
object's subjective determinations (qualities and so forth) and its essence, as
this would come forward before subjects of another kind than we?

Here again we meet the irrational, and here we see perhaps most clearly how
inexhaustible existence is in relation to our knowledge. What was justified in
Kant's setting up of the concept ``Ding an sich'' lay in this, that a
limit-concept is needed to express the irrational relation between what he called
the ``form'' of our knowledge and its ``matter''. If we would keep the concept
``thing in itself'', we can employ it in Kant's spirit and yet avoid falling into
the contradictions with which it is burdened in his philosophy, if we use it to
express the fact that the difference between subject and object always comes into
force anew, however often we may have found an objective explanation of the
subject's peculiarities (of \textit{S}\textsubscript{1} by
\textit{O}\textsubscript{1}) and a subjective explanation of the object's
peculiarities (of \textit{O}\textsubscript{1} by \textit{S}\textsubscript{2}).
The irrational shows itself in this, that a continual series-formation (of the
type: \textit{S}\textsubscript{1}\,\{\,\textit{O}\textsubscript{1}\,\{\,%
\textit{S}\textsubscript{2}\,\{\,\textit{O}\textsubscript{2}\dots) is possible
and necessary. Thought must ever anew set about finding predicates for the
determination of existence, because the source that makes the stream of thought
possible is inexhaustible. ``The thing in itself'' is the obscure, tacitly
understood initial representation, which comes forward ever in new ways and
demands new determinations.\footnote{Compare \textit{Det psykologiske Grundlag
for logiske Domme.} p.~39\,f.; 46\,f., and (as regards religious judgments)
\textit{Religionsfilosofi}, p.~67--71.} Perhaps the right symbol for the relation
of our knowledge to existence would not be an irrational number, but an imaginary
number, since existence might have attributes which cannot at all be reached or
determined in the dimensions in which our thought is able to move. The
possibility of this cannot in any case be disproved. Nor can that other
possibility, that existence is rational only in a very slight degree, and that it
might one day turn towards us another side, one that would not make possible the
building of coherent and applicable series of thoughts; there would then, as I
have expressed it in another place,\footnote{See my treatise \textit{Filosofien
og Livet.} (Tilskueren 1901). p.~34. --- On the imaginary in an epistemological
sense compare \textsc{Wundt}: \textit{System der Philosophie.} p.~195--199.}
begin a logical ice age for us. The relation between subject and object would not
come into force at all; neither an \textit{S}\textsubscript{\textit{o}} nor an
\textit{O}\textsubscript{\textit{s}} could subsist.

When in an earlier connection I used Schiller's words: Eng ist die Welt, und das
Gehirn ist weit, it now shows itself that the sentence can be reversed: Gross ist
die Welt, und das Gehirn ist klein! Knowledge is, however rich and powerful it
may be, always only a part of existence, and the problem of knowledge could be
solved only if existence as a whole (in so far as it can be thought as a
completed whole!) could be expressed by the help of a single part of it. Symbolic
the expression must in any case always remain. Our knowledge gives us, even where
it reaches highest, only a section of something more comprehensive. Of all the
possibilities of thought only a single one is represented in the reality we know;
but the reality we know is itself always only a part of a greater whole --- and
the relation between part and whole we cannot here determine. We can form no
exhaustive concept of reality.

% ============================================================
%  III.  TILVÆRELSESPROBLEMET  (pp. 52--69; notes 52--56)
% ============================================================
% Fisher renders this chapter "The Problem of Being"; "Tilværelse" is kept as
% "existence" here for consistency with the Introduction, where Høffding names
% the problem "Problemet om Tilværelsens Natur" and pairs it with the valuing
% personality's existing. Fisher's section heads: 1. Problem and method /
% 2. Metaphysic as an art / 3. First type-phenomenon (unity or plurality?) /
% 4. Second type-phenomenon (spirit or matter?) / 5. (rest or development?).
\fpchapter{III.\quad The Problem of Existence}
\label{ch:tilvaerelse}

\noindent\textbf{1.}~Both the problem of consciousness and the problem of
knowledge point out beyond themselves: the life of consciousness is a part of
existence, and it is the task of knowledge to understand existence. It is then
natural to ask what place the life of consciousness occupies in existence, and
what picture knowledge is able to give us of existence. The problem to which we
are thus led over we may call the \emph{cosmological}, since it is the problem
whether a completed world-view is possible: \emph{Kosmos} means existence
regarded as a whole, and \emph{Logos} means doctrine, view. The word metaphysics
it will also be warranted to use here, though it has been employed in many
different senses and is on that account less clear.

The chief difficulty in this problem will be that the cosmological principles,
the principles which are to lie at the base of the world-view, cannot without
more ado be fetched from any special domain of experience, since the point would
precisely be to connect all domains of experience into a whole and to give a
positive characterization of it.

An attempt to treat the problem of existence will have both a formal and a real
character: partly it must be asked what demands must be made on the concept that
is to be able to embrace existence as a whole, and whether it is possible to
satisfy those demands; partly it must be asked what positive marks can be
assigned of such a whole, if we are able to form a concept of it.

The \emph{formal} motives for speculation on the nature of existence lie deep in
the constitution of consciousness and of knowledge. They hang together with the
need for continuity, the need in which personality and science meet. The very
essence of thinking expresses itself on all levels and in all forms as a
gathering-together, a synthesis, and it is then not to be expected that thinking
should of its own accord halt in its attempt to work together what is
sporadically given. This striving acquires a special and practical significance
from the fact that a firm and continuous connection is the only criterion we have
in doubtful cases when it is a question of distinguishing between dream or
illusion and reality. The more comprehensive and the more inwardly coherent a
concept of reality we can form, the greater the certainty it contains. There will
therefore, on theoretical as well as practical grounds, be an inclination to go
to the limit and to attempt to continue and to complete what theoretical and
practical orientation already leads us into. The ideal would be reached if a
complete harmony could be exhibited between all our experiences, a continuous
whole into which all the special domains of experience join themselves according
to their own laws.

But the discussion of the problem of knowledge must already have taught us that
such a completed world-view is impossible, and that in a certain sense it
contains a self-contradiction. None of the special domains of experience lies
before us completed; new experiences and new riddles are always coming up; our
comprehending thought is always getting new tasks. And since our knowledge always
works by holding things together and comparing them, every total picture, if it
were to be an object of completed knowledge, would have to be held together and
set beside something that was different from it: only so could it get its full
determinateness; but if there were anything different from it, it would after all
not be a \emph{total} picture! --- It is here indifferent whether we stop at the
totality we have attempted to form, or whether we go back to the principle of
such a totality and call it ``God'' in contrast to the totality itself as ``the
world''; the antinomy remains the same in both cases.\footnote{Compare on this my
\textit{Religionsfilosofi}, p.~57--65. --- In the epistemological section of the
\textit{Religionsfilosofi} I have on the whole already set forth my conception of
the problem of existence and of its treatment.} The irrational announces itself
here again, as it did at the problem of knowledge. And there is here surely an
inner connection. That knowledge could not be completed might hang together with
existence itself not being completed, not being finished, but being in continual
becoming, as the single personality is, and as knowledge is. Perhaps it also
conceals simultaneous disharmonies within itself, which make it impossible for it
to constitute a harmonious whole. The analogy between the different problems
would in that case be especially clear. In the various systems men have been too
sure that existence in itself was a closed and constant whole, and that it was
only our will and our thought that had always to struggle in order to subsist and
to win harmony.

A \emph{real} motive for speculation lies in the prominent role which a
phenomenon, or a domain of experience, or a side of existence can play for us.
From this there can arise a tendency to lay this phenomenon or this side at the
base of an interpretation of existence, and to attempt to derive the other
phenomena or sides from it, or in any case to trace them back to it. The
different cosmological systems are so many attempts to sound the depths of
existence, to test how far a thought can cast light over existence. They are a
series of thought-experiments by which the bearing power of our most important
thoughts, their capacity to serve as the foundation of a comprehensive
world-view, has been tried.

Every such attempt goes --- as \textsc{Leibniz} first saw with brilliant clarity
--- necessarily by the way of \emph{analogy}. --- In all science analogy can
become of importance. All designations for psychical phenomena were originally
formed on the basis of an analogy with physical phenomena, and in psychology we
are constantly compelled to operate with physical analogies; it then becomes an
important question how far their validity reaches. The causal principle
expresses, according to Kant, that between the actual events in the world there
is a relation analogous to that between ground and consequent in our thinking.
The atomic theory, and on the whole the so-called mechanical view of nature, is
in the most recent time regarded epistemologically as a great analogy, by whose
help the qualitative changes in nature can be described and calculated. In
discoveries too analogy plays a great role: the analogy between chemistry and
physics helped Robert Mayer to his discovery of the conservation of energy. But
when analogy is employed metaphysically or cosmologically, it is not \emph{one
particular} domain of existence that serves to illuminate \emph{another
particular one}; it is \emph{one particular} domain that is used to express
existence \emph{as a whole}. This symbolizing becomes of another kind and another
validity than the symbolizing which finds application within the special domains:
it cannot be carried through with full consistency, and it cannot be verified.
Neither the warrant of the analogy nor its limitation can be strictly exhibited.
In this the cosmological or metaphysical symbols differ from the scientific ones.
--- As I have sought to show in my \textit{Religionsfilosofi}, the religious
symbols share the fate of the metaphysical. In both cases an attempt is made to
form concluding concepts with real validity; the difference rests only on the
motive that leads to these attempts.

In the problem of knowledge we came likewise to an irrational relation between
part and whole. There it was a question \emph{how far the whole could be
expressed in and for a single part}. Here, at the cosmological problem, it is a
question \emph{how far determinations can be fetched from a single part which
hold for the whole as such}. It is one and the same relation, coming forward in
the two problems from different sides. ---

The character and the value of a world-view depend not only on how clearly and
consistently the analogy is carried through, but also on where the analogy is
fetched from. The phenomenon (the part or side of experience) on which the
analogy is built we may call the \emph{primal phenomenon}. This expression stems
from \textsc{Goethe}. But Goethe understood by a primal phenomenon not merely a
phenomenon of a typical kind, able to illuminate other phenomena for us; it means
for him a fact which is at the same time a law, so that one need only utter it in
order to have explained it, and it must not be regarded as composite, since it is
itself to be the symbol of everything else, to be a ``limit of our beholding''
which at once awakens ``the great wonder'' and the feeling that we stand at the
limit of our capacity!\footnote{\textsc{Goethe}: \textit{Farbenlehre.} §~124;
175--177. Compare \textit{Gespräche mit Eckermann.} 18. Febr. 1829; 21. Decbr.
1831. --- Goethe applied this concept not only in the domain of the theory of
colours, but in magnetism and in botany as well. See \textit{Sprüche in Prosa.}
(Ueber Naturwissenschaft) and \textit{Gespr. mit Eckermann.} 27. Januar 1830. ---
\textsc{Hegel} was (as is seen from Goethe's \textit{Nachträge zur Farbenlehre})
much taken with the concept of the primal phenomenon as Goethe had set it up.}
--- It is well known what a fateful role this concept played in Goethe's theory
of colours. By his very definition of it he meant to exclude every further
investigation and explanation of the phenomenon he chose as ``primal
phenomenon''. But this defect need not cleave to it. It may very well designate
the point in our experience from which we attempt to orient ourselves in all
directions, and yet be itself an object of further investigation and explanation.
The main thing is that the phenomenon in question should be so peculiar and so
significant that it can take on a typical character for us. The interpretation of
existence must always proceed from a single domain, and it can proceed without
that domain being itself withdrawn from special scientific treatment.

But where now are we to seek the primal phenomenon? On that turns the battle
between the different world-views. Now it is life, now thought, now matter that
is made the primal phenomenon and laid at the base of the interpretation. Before
we go over to considering the most important primal phenomena that have been
taken into the service of cosmological interpretation, we must nevertheless dwell
a little on this interpretation itself.

\noindent\textbf{2.}~Metaphysics can become dogmatic, and thereby work harm in a
scientific respect, in two ways --- partly by cutting off every special
scientific explanation of the chosen ``primal phenomenon'', partly by conceiving
its analogizing method as more scientific than it is, and forgetting the
necessity of continual confirmation by experience. It cannot interfere in the
housekeeping of the special sciences, in that of the human sciences as little as
in that of the natural sciences. Its philosophical place is at the boundary of
scientific thinking; its task is to give a last interpretation. The theory of
knowledge led us to a limit-concept for that which conditions the continual
conflict in our conception of the world between quality and quantity, the
elementary and the ideal concept of cause, subject and object. If we choose to
call that which contains the ground of these oppositions ``the thing in itself'',
then this expression designates the philosophical place of metaphysics or
cosmology.

The last interpretation will have more of the stamp of art than of science. The
transition from the empirical and critical parts of philosophy (psychology and
the theory of knowledge) to cosmology presents a certain analogy to the
transition from historical criticism to the writing of history. The historian
forms the critically treated material into a total picture of events and
characters; what is fragmentary is filled out, and a rounding-off is given in
which the personality of the one who presents it necessarily plays its part. The
cosmological philosopher will in the same way give us a total picture of
existence on the basis of our knowledge of it; he will work the scattered
features together into a totality within which a single element (the primal
phenomenon) is particularly determinative, so that the timbre of the whole
depends on it. In the choice of the primal phenomenon and in the carrying through
of the analogy a decidedly personal element asserts itself. A great philosophical
system is a work of art, a drama in which the single speeches and scenes join
together, now more with logical, now more with aesthetic-ethical harmony. The
systems of Aristotle, Spinoza and Schopenhauer present this character in a
special degree.

This art will become harder and harder to practise. For in future it will
presuppose not only that a rich material stands at one's disposal and can be
shaped with constructive power, but must presuppose as well the ability to avoid
dogmatizing, to maintain the significance of ideas and of constructions of
thought without confusing them with absolute truths. The thing is, as
\textsc{Lessing} said, to think gymnastically, not dogmatically. The art will be
this: to unite bold development of thought with a wakeful critical consciousness.
It is this side of the matter that I have particularly emphasized in my treatise
\textit{Filosofi som Kunst}. Earlier the artistic element in philosophical
thinking was emphasized by \textsc{Schopenhauer}, who however thinks mostly of
the \emph{involuntary} coming-to-be of the decisive construction of thought,
whereby art's ``What?'' takes the place of science's ``Why?'' When
\textsc{Albert Lange} too called philosophical speculation an art, he was
thinking rather of the \emph{idealizing} tendency, the need to see in ideal
images an expression of the highest reality.\footnote{See my \textit{Mindre
Arbejder}, p.~180--194. --- On Schopenhauer and Lange see \textit{Den nyere
Filosofis Historie.} II. p.~198\,f.; 501\,f.} --- At the religious problem we
shall meet a relation akin to the transition from science to art of which we have
here spoken; only that there it will be a question of a completion of the view of
\emph{life}, whereas here it is a question of a completion of the view of
\emph{the world}. ---

It is not, as the history of philosophy shows, at all times that the conditions
are present for the exercise of the art here spoken of in the grand style. A long
period of accumulation of material and of points of view, a great concentration
of mind, and a thought-energy awakened by the sharpening of the problems must be
present. We find in the history of philosophy that the systems lie gathered in
certain single periods, and even bound to single places as well. The fifth and
fourth centuries before Christ created the great Greek systems, and here Athens
was the philosophical centre. In the seventeenth century arose the principal
systems of the modern age, which came into being on the basis of the manifold
material and the new points of view of the Renaissance and of the newly founded
natural science; here Holland occupies a central position. At the beginning of
the nineteenth century all the systems of philosophical idealism were founded in
Jena, where Kant's disciples met with Goethe's disciples and with the champions
of Romanticism. It is at such points of time that philosophy stands as one of the
symptoms of the course of the life of the spirit; in the intervening times the
more special discussion of psychological, epistemological and ethical problems
goes its quieter way. ---

I now go over to bringing forward the most important primal phenomena which can
serve as a systematic foundation and have historically rendered such service.

\noindent\textbf{3.}~The first fact which it lies near for philosophy to regard
as the primal phenomenon in existence is this, that existence is to a great
extent \emph{intelligible} to us. We can recognize phenomena, infer from one
phenomenon to another, and find continuity between them. Even though our
hypotheses are really working hypotheses, and even though the old, naive concept
of truth must more and more give place to the dynamic or symbolic concept of
truth (see II, 2), the very fact that with our faculties and our methods we can
to a certain degree orient ourselves in existence must nevertheless hang together
with the essence of existence. The applicability of a method always contains a
testimony about the constitution of its subject-matter. The intelligibility of
existence for us points to an inner unity in its essence, which lays itself open
to the day in the lawfulness that the course of phenomena presents. Our very
criterion of reality consists in firm connection, and it is a natural extension
of what the application of this criterion establishes when we find in existence a
\emph{unifying power} by which the single elements and events are bound together.
Philosophy has a natural tendency to go back to such a principle; Plato and
Giordano Bruno, Spinoza and Hegel, Fechner and Lotze worked in this direction,
and even Kant, for all his great critical scruple, has borne witness to the
significance of this thought. While popular thinking is inclined to seek a
completion by climbing the ladder of causes so as to end in a first cause --- a
road that leads only to an endless series --- philosophical thinking seeks its
completion rather in depth than in length, bethinking itself of what the
presupposition is for a rational world-picture, however hypothetical, being
formed at all. \emph{Causality} stands here as the primal phenomenon, and it is
not least for this reason that the discussion of the causal principle fills so
large a place in the history of modern philosophy.

Against the monism which thus presses forward one might object that what is
characteristic of the scientific understanding of phenomena is not only the
assumption of an inner connection between them, but also the assumption of a
manifold of elements between which the interaction takes place that the laws
found by science hold for. Why not regard this manifold as the primal phenomenon,
so that our metaphysics becomes pluralistic and not monistic?

The answer to this question is given in the fact that all the properties or
forces with which we equip the elements of existence (whether we represent them
to ourselves as material atoms or as psychical monads) are fetched from the
lawlike connection in which those elements are members. The properties of things
are the constant ways in which they affect one another and are affected by one
another. ``Force'' is here the element which for us contains the ground of a
change in one or more other elements. So it goes with the kindred concepts
energy, capacity, individuality. All these concepts are derivative in relation to
the concept of law. In the case of individuality one must here think not only of
the law that holds for the individual's relation to the surrounding world, but
also of the law for the relation between the individual's inner states.

In the same direction points another consideration. All knowledge begins as
analysis of given intuitions (observations or recollections). If such an
intuition is to be able to become an object of analysis, it must come forward as
a totality embracing an aggregate of elements (parts or properties). Every act of
knowledge rests on a section cut out of a greater whole: every conceptual
determination (definition) is a limitation; every judgment rests on
presuppositions that point out beyond the act of judging itself; every inference
presupposes several premisses whose validity must often be grounded along highly
different roads. We always move within a more comprehensive whole, in which the
conditions must be sought for the single results we arrive at.

When thus, both in the world of reality and in the world of pure thought, the
particular gets its constitution and its validity from the unitary connection in
which it occurs, monism will stand as taking a more fundamental starting-point
than pluralism.

Meanwhile the theory of knowledge does show that there are limits to
understanding. The empiricist and the sceptic will always be able to stop the
monistic metaphysician by pointing to the actual limits of knowledge. Not even
the criterion of reality can we apply everywhere, and so we cannot carry through
the strict difference between dream and reality. With the same right with which
one infers from the possibility of rational knowledge to a unifying power in
existence, one would then, it seems, be able to infer to an irrational element in
existence, to a cosmological principle which prevented the elements of existence
from standing in a rationally determinable relation to one another.

From the monistic side it may again be answered that if the unifying power does
not rule everywhere, that might indicate that existence itself is in becoming, in
development, and that what stands for us as a lawlike connection presupposes a
development in the innermost of existence which is not yet completed. It would
then, from this point of view too, be \emph{time} that conditions the irrational.
And so long as thought, knowledge, is not finished but is caught up in becoming,
so long \emph{can} existence as a whole not be finished either. Thought,
knowledge, is after all a part of existence! And when \emph{it} is in becoming,
there might well be more that is in becoming. It is the strange contradiction of
the great rationalistic systems that even if they could explain everything else,
they could not explain the striving and labouring thought whose products they
are. In Plato and Spinoza, in Hegel and Boström this contradiction comes forward.

\emph{Critical monism}, which holds fast to the validity of time and thereby to
the incompletability of existence as well as of knowledge, may very well build
its world-view on causality and rationality as primal phenomena. It then finds in
existence a struggling unifying power which through progressive development is
carried out beyond what is sporadic and conflicting. Perhaps new elements are
always arising which have in their turn to be worked together, and development
must then always begin anew. The work of thought is hereby seen in a cosmic
light: a constant connection between our methods and hypotheses and the real
processes of existence is after all the goal thought sets itself (even when it
exchanges the static concept of truth for the dynamic); and if thought succeeds
in approaching this goal, existence will thereby have become more rational than
it was before, since a new constant and harmonious relation has been realized.
The thinker to whom it is given to work forward in this direction may with right
say: ``Ort für Ort \textit{sind} wir im Innern'' --- however distant and exalted
thought's highest ideal may stand before him.

\noindent\textbf{4.}~Can we not attempt a real, positive determination of the
principle of unity which, on the hypothesis just developed, holds existence
together in its innermost? --- Every attempt in this direction must make use of
analogy in a particularly high degree. We have here no such fundamental fact to
point to as we had in setting up that hypothesis, where we could support
ourselves on the connection of the problem of existence with the problem of
knowledge. In a closer determination of the principle of existence, thought will
be led to the most fundamental difference which the phenomena accessible to us
present: the difference between the spiritual and the material. There seems here
to be a forced choice for every attempt in a cosmological direction.

How I stand, in a purely methodical and empirical way, to the problem of the
relation between spirit and matter I have already developed (I, 3). But the
question comes up here again from another point of view. For whether we are
methodically and empirically dualists, materialists or monists, it may be asked
how we think the fundamental essence of existence --- what the underlying
attribute is, that under which we finally --- when we think our last thought ---
think existence.

If I lay at the base the analogy with what is extended and movable in space, then
my metaphysics, and not merely my method, becomes materialistic. As a concluding
hypothesis \emph{materialism} has the advantage that it gives us the picture of a
great continuity, and that we need not leave what is intuitable. Materialism is a
childlike and naive conception; it is man's first philosophy. But the impression
of the connection and the extent of the material world exercises a constantly
overwhelming power, so that attempts in a materialistic direction come forward
again and again, even if since the appearance of the critical philosophy they
cannot present themselves with the dogmatic assurance of former times. They will
always run aground, partly on the impossibility of tracing the spiritual back to
the material, partly on the epistemological scruple that we have the material
always only \emph{as object} for consciousness, and that if materialism were
right there could exist nothing \emph{for which} the material could become object
or phenomenon. \textsc{Hobbes}, whose thinking went strongly in a materialistic
direction, nevertheless halted at precisely this consideration, that of all
phenomena \emph{the very fact that something can become a phenomenon for us} is
the most important.

In contrast to materialism, \emph{metaphysical idealism} lays at the base the
analogy with spiritual phenomena. It can empirically espouse dualism (like
\textsc{Lotze}) or materialism (like \textsc{Schopenhauer}) or Spinozistic monism
(like \textsc{Leibniz, Fechner} and \textsc{Wundt}). The last thought is in all
forms this, that only the analogy with the spiritual states we know from
ourselves gives us a key to the understanding of the innermost essence of
existence. Ourselves alone do we know from within, everything else merely from
without! Psychological experience already teaches us that there are many degrees
and kinds of spiritual existence, and if we would use it for the interpretation
of the whole of existence, of its innermost ground and its ultimate elements,
then we must naturally think these series of degrees and of qualities continued
without limit. This idealistic cosmology can be varied in many ways, and appears
historically in many forms, conditioned especially by the other motives that
determine the world-view (for example the tendency to optimism or to pessimism,
chief weight on thought or on will, and so forth). It appears already in the
Indian ``Upanishads'': in the doctrine that Bráhman (the world-principle) is
Atman (soul).

Sometimes it is denied that it is an inference from analogy that one here relies
on, and it is maintained that metaphysical idealism (or idealistic cosmology) has
been reached along the strict road of logical construction, of the dialectical
method. But even the proudest edifice of thought that men have believed
themselves able to raise in this way, \textsc{Hegel}'s system, really comes to
showing that everything in existence hangs together as inwardly as the thoughts
in the human mind: it is in reality the latter that is used as the basis for an
analogizing, since it stands as the best example of an inward
totality.\footnote{See \textit{Den nyere Filosofis Historie.} II. p.~165\,f. ---
It is characteristic that the English thinkers who in the most recent time have
worked with some of Hegel's fundamental thoughts concede that what is in question
here is analogy rather than demonstration. \textsc{Francis Bradley} does not
consider even the analogy defensible; the highest reality can, according to him,
be called neither soul nor body, though the soul more than the body has that
combination of extent and unity, that ``self-consisting'', which is the mark of
true reality. (\textit{Appearance and Reality.} p.~307; 357). --- In his treatise
\textit{Hegel's treatment of the categories of the idea} (Mind 1900. p.~149\,ff.)
\textsc{Mc.~Taggart} states that we have only a single example of the category
``spirit'', which in Hegel acquires a theological or cosmological significance.
--- When one distinguishes (as \textsc{A. E. Taylor} in Mind 1900 p.~245) between
two idealistic schools, one standing nearest to Leibniz and Lotze and one
standing nearest to Hegel, I must maintain that it is the first of these schools
that has most clearly seen the real philosophical foundation of metaphysical
idealism. (That it is an analogy with the human mind that lies at the base in
both cases, Taylor really concedes himself when he says that both tendencies are
agreed on ``the main principle that it is in mind, and nowhere else, that we are
face to face with the central reality of the universe''. ibid.)} ---

The idealistic analogizing could lead to a concluding solution only if we had
means of determining in a positive way the different degrees and kinds of
spiritual existence which must be found in the world if metaphysical idealism is
right. But we cannot, as has already shown itself in an earlier connection
(I, 3), get further here than to the negative concept of potential psychical
energy. A verification of idealism is therefore impossible. And even if it could
be carried through in full, there would still remain the difficulty that the
material can as little be derived from the spiritual as the spiritual from the
material.

But it is by no means certain that we have a forced choice between a
materialistic and an idealistic solution of the problem of existence. The
difference between matter and spirit is certainly a chief difference in the
content of our experience; but no proof can be given that there are not in
existence other fundamental forms (``attributes'') than these two. \emph{If} the
problem of existence is to be solved out of human experience, then one of those
two possibilities must be chosen. But the question is whether our experience
presents us with elements enough for a real solution. The richness of existence
may be much greater than the possibilities of our experience. Here again it may
hold that the world is great and our brain is small; here again we strike on the
irrational. Could it be proved that the difference between spirit and matter was
contradictory and not contrary, so that an absolute either-or lay before us, the
problem of existence would be simpler than it is --- and yet more involved than
human thought, which is inclined to think itself fully equipped for religious and
metaphysical speculations, has often supposed. \emph{Critical monism}, which
strives to hold fast the thought of unity without dogmatizing, must impress upon
us that we here lack the conditions for a complete solution. That there is a
possibility of more forms than our experience presents may indicate that the
whole problem lies deeper than has been supposed. There might, for example, be a
fundamental property of existence of which both spirit and matter were
consequences, and the insolubility of the problem for us might come of our not
knowing this fundamental property.

\noindent\textbf{5.}~A third primal phenomenon we have when we have made our
choice between subsistence and development (being and becoming). Ancient thinking
was throughout inclined to hold to subsisting being; it was a philosophy of
concepts, which above all sought to trace phenomena back to fixed concepts of
kinds. Plato's theory of Ideas forms here the high point (compare II, 1). And it
--- or the same tendency of mind that has found expression in it --- comes
forward not seldom in more recent inquiry as well, and that not only in
philosophy but everywhere in science where the battle is fought for fixed,
unchangeable species or types. Even where ancient thinking assumed a development,
it was particularly inclined to assume a rhythmical process which brought the
same states and events with it again and again. Development as an ever
forward-advancing series of changes is a modern idea, due to widened experience
in the domain of history as well as of nature, but one which has also taken shape
under the influence of the Persian-Christian type of religion.\footnote{Compare on
this type my \textit{Religionsfilosofi.} p.~50; 121\,f.; 273\,f.}

That the concept of development comes to play so great a role for the thinking of
the modern age hangs together with the causal concept's having come forward so
strongly. The more the elementary concept of cause, through continued inquiry,
approaches the ideal concept of cause, the more does the causal relation come to
designate a continuous process in which the following members are determined by
the preceding. And the concept of development comes forward as soon as it shows
itself that the \emph{direction} is essential, so that time cannot be reversed
(compare II, 4). While the formal sciences, which rest on the principle of
identity, can go forwards and backwards in their series of thoughts, the
time-relation, and particularly the direction of time, acquires a decisive
significance for the real sciences, and these thereby retain a historical
character. Here the very concept of event is a primal phenomenon. But in the
concept of development there lies this, that not merely does something happen,
but that through the series of events states are reached which have a certain
character of wholeness, a manifold of elements being connected in such a way
that, in spite of their difference of kind, they work in firm connection with one
another and with a certain completeness over against their surroundings. Most
fully, and with a great wealth of examples, \textsc{Herbert Spencer} has analysed
the concept of development and maintained the combination of differentiation and
concentration as its essential feature. It is by this concept as measure that we
determine what --- from a purely theoretical point of view --- is to be called
``lower'' and ``higher'' within a changing series of states or forms.

In spite of its connection with the causal concept, the concept of development is
an independent concept which cannot be derived from the general concept of cause
(though it presupposes it). The causal principle would hold even if only a
rhythmical oscillation took place, without progressive formation of new states of
wholeness. Development stands as a fact of experience which casts light on the
nature of existence. It was in and for itself quite possible that the different
causal series in existence should either not meet at all, or else meet only in
order to collide in a disharmonious way. Now experience shows that under certain
conditions they can meet in such a way that they are combined into more composite
processes and give rise to peculiar wholes. The coming-to-be of the stellar
systems, the arising, the subsistence and the unfolding of life, the existence of
the life of the soul and of social, historical human life are testimonies to an
individualizing and totalizing tendency in existence. These phenomena constantly
set inquiry the greatest problems; when we point to them here as primal
phenomena, we decide nothing about whether these problems can be solved or not,
but regard them merely as characteristics of existence as it once for all is. It
is certainly --- this must be added here --- a great misunderstanding when it is
often supposed that these phenomena would be more remarkable if they could
\emph{not} be explained by the help of laws that science can find. On the
contrary: if a ``natural'' explanation of them can be found, it will be possible
to infer with greater certainty than before that the individualizing and
totalizing tendency lies deeply grounded in the kernel of existence.

But how deeply? Experience shows us not only, and often, a purely external and
indifferent relation between different causal series, but also, and often,
collisions between them, or at any rate disharmonious relations which hamper or
hinder the arising or the subsistence of individualities and totalities. And
after development follows dissolution, and the question becomes whether through
the rhythmical alternation of these processes a progressive course of development
can be traced, or whether we must halt at the ancient thought of rhythm as the
last thought.

Here the irrational comes forward again, and it comes forward here in a sharper
form than at the earlier points where we met it; it designates here what hampers
and dissolves, the tendency to remain standing at, or to return to, the most
elementary forms of existence. Existence gives us here, the more inquiry
penetrates into it, the picture of a struggle --- the great struggle which all
forms of existence that bear the stamp of individuality and of totality must wage
in order to subsist. The struggle itself is again two-sided: it can be a means to
development, but it can also be a road to destruction; and which of these
possibilities is the preponderant one?

At the same time there shows itself here still more clearly than at the earlier
points the impossibility of forming an absolutely completed concept of existence
as a whole: if the conflict between elements and totalities is an
\emph{essential} feature of existence, then it is a feature that can always be
found only in limited sections of existence, not in existence regarded as a
whole, since an absolute whole can meet no outward resistance, can be engaged in
no struggle. We must content ourselves with saying that wherever we come upon
existence, we find that within it such a conflict between elements and totalities
goes on. The great question is whether through this conflict it is the elements
or the totalities (the solar systems, the organisms, the souls, the societies)
that carry off the victory. Empirically we stand here in the midst of the great
world-process, and can get no further than to maintain that it is no purely
subjective measure we apply when we designate one form of existence as higher or
lower than another, since development is a phenomenon characteristic of all the
existence we know --- a primal phenomenon. ---

By the result at which we thus end in the matter of the problem of existence, a
natural transition to the last problem, the problem of valuation, becomes
possible.

If existence were finished, harmonious and unchangeable, no ethics would be
possible. All ethics demands a work performed; but there would be no room for any
work if everything consisted in eternal and actual completedness. Precisely this,
that all individualities and totalities must struggle in order to subsist ---
that there are constantly disharmonies to overcome, sporadic tendencies to
connect --- precisely this makes an ethics possible. For ethics investigates the
principles of a carrying-further of the individualizing and totalizing tendency
which experience exhibits in existence. Ethics rests on a self-contradiction if
the course of the world does not partly proceed, or cannot partly proceed,
through human will, just as it partly proceeds or can proceed through human
thought-work (III, 5). Ethics then takes up the problem of continuity where the
first three problems let it go.

Still more clearly does the religious problem find a point of attachment in the
foregoing consideration, since it --- as I shall attempt to show --- concerns the
continuity of values under the struggle of existence.

% ============================================================
%  IV.  VURDERINGSPROBLEMET  (pp. 70--84; notes 57--63)
% ============================================================
% The numbered run-in divisions here run 1--7 STRAIGHT THROUGH and do not
% restart at B: §1 is a chapter introduction standing before A, §§2--4 fall
% under A (a = §2, b = §3--4), and §§5--7 under B.  Verified against the epub's
% own h2/h3 headings.
\fpchapter{IV.\quad The Problem of Valuation}
\label{ch:vurdering}

\noindent\textbf{1.}~That the problem of valuation is set up as a separate
problem presupposes a distinction between understanding and valuation which has
not always been acknowledged, and which has not yet made its way through. There
has been an inclination in philosophy to let the two run into one: thus Plato's
``Ideas'' and Spinoza's ``substance'' are at once an expression of the two
thinkers' understanding of existence and of their admiration for it. In any case
men have been inclined to let valuation be a mere consequence of understanding,
while conversely mystical and theological tendencies have been inclined to let
understanding be dependent on valuation. Wholly independent of one another they
certainly are not: understanding has its value, and would not be possible if that
were not so, since men would not otherwise strive after it --- for some, indeed,
the joy of knowledge is the highest of all; and valuation can become an object of
psychological and historical understanding, since all value rests on the relation
of experiences and states to life at its various levels, to life's subsistence
and its development. The problem of valuation will therefore also show itself to
present an analogy to the three first, theoretical problems. It is in the domain
of valuation, as in the domain of personality, of knowledge and of existence, the
principle of continuity that leads to the decisive problems.

\emph{Value} belongs to everything that brings a satisfaction with it or
relieves a need. Sometimes it is only through a satisfaction's arising that we
become aware that there has been a gap in our existence; sometimes this gap is
felt beforehand and causes a want or calls forth impulse and desire. When what is
valuable does not immediately fall to our lot, we make it an \emph{end} and seek
means of attaining it. Whatever stands as a means of attaining something
immediately valuable acquires a mediate value for us. In our values and our ends
the innermost essence of our feeling and our will makes itself known. Just as the
concept of end depends on the concept of value, so the concept of \emph{norm}
depends in turn on the concept of end: the norm is the rule for the activity
necessary to the attainment of the end.\footnote{In my \textit{Etik} (2nd ed.,
p.~21) I did not enter more closely into the relation between the concepts value,
end and norm. See instead \textit{Det psykologiske Grundlag for logiske Domme.}
p.~47 (compare p.~65) and my \textit{Religionsfilosofi.} p.~12.} --- It had a
fateful influence on the treatment of the problem of valuation that
\textsc{Immanuel Kant} wanted to reverse the relation and to derive the concepts
of end and of value from the concept of norm (``law''). That is a psychological
impossibility.

Now experience shows that for different individuals, and for the same individual
at different times, different values hold. I have already mentioned the value of
knowledge; besides it there are to be mentioned especially the values attached to
the impulse of self-preservation and to the subsistence and development of life
in larger or smaller circles --- to movement and activity --- to images (those of
reality or of fantasy), and so forth. If a comparison is to be instituted between
different values --- and every conscious valuation consists in such a comparison
--- then a \emph{fundamental value} must be presupposed, from which the order of
rank of the different values can be fixed. A determinate value must be laid at
the base and furnish the measure for all other values, if consistent thinking and
discussion is to become possible in the domain of values. From a given
fundamental value there can then --- on the presupposition of sufficient
experience --- be constructed a system of valuation in which every single value
gets its place according to its relation to the fundamental value. Thinking,
``practical reason'', is needed in order to determine this relation, and also in
order to find means and ways of producing or finding out the several values under
certain given conditions. But on bare ground no valuation (or ``revaluation'') is
possible. The concept of fundamental value is, within the problem of valuation
(that is, in ethics and the philosophy of religion), what the concept of primal
phenomenon is within the problem of existence (that is, in metaphysics). ---

The problem of valuation divides into two problems, the ethical and the
religious. Ethical valuation concerns human actions, qualities and orderings of
life (institutions); religious valuation reaches further and values existence
according to the fate of values in the world of reality. Within both problems the
relation between continuity and discontinuity will --- as at the earlier problems
--- show itself to be of decisive significance.

\section*{A.\quad The Ethical Problem}
\label{sec:etiske}

\subsection*{a.\quad Ethical Work}
\label{sec:etiske-arbejde}

\noindent\textbf{2.}~There is, as the discussion of the problem of existence has
shown, room for a work by which existence is developed further. Such a work is
performed in all human culture, but quite particularly in ethical striving. This
may be designated as a striving to bring forth greater continuity, partly within
the single personality, partly between the different personalities among
themselves. The measure of ethical striving, and so the principle of ethical
valuation, will show itself to be determined by the principle of continuity.

Personality presupposes continuity, and this is in turn conditioned by there
being one fundamental value which determines the value of the several moments,
periods of life, capacities and impulses. Development into a true personality
presupposes a striving out beyond what is momentary and what is sporadic --- an
overcoming of the tendency of the single moment or the single need to isolation
and sole dominion. The thing is to work the several moments and elements
harmoniously into the life of personality as a whole. Here is a work to be done,
a struggle to be come through, which in different individuals demands different
degrees of energy.

There are ethical standpoints which assert the right of the several moments and
impulses, of the successive and simultaneous differences, over against the whole
of life. But this does not conflict with the measure's being fetched from the
principle of continuity. Such standpoints may be warranted in so far as they
demand that the subordination of moments and special impulses to the whole of
life be grounded. Continuity does not mean absence of difference, but the
ordering of the differences into a graded series. Life as a whole can always be
called to account by the single element of life. It will always stand as an
imperfection when a moment, a period of life, a capacity, an impulse is made a
\emph{mere} means without acquiring independent value. The art will be to let
them have immediate and mediate value at once. --- And in so far as the ethics of
the moment or of the speciality appears in absolute form and proclaims the
sovereignty of the single moment or element, one will often find at the base the
presupposition that life as a whole is represented by, or is concentrated in, the
single moment or the single need, so that all other considerations pale. So it is
in the great moments of self-sacrifice, where a man --- as \textsc{Aristotle}
says --- ``would rather perform one great and beautiful action than many small
ones''; and so it is where a man sees the calling of his whole life in the
special development of a single capacity. Or one believes in a harmony between
all the different moments of life, so that a complete absorption in each of them
becomes possible without the satisfaction of the single moment depriving the
other moments of anything. So \textsc{Aristippus} seems to have thought of it. Or
the highest is placed in a hovering above the several moments and occupations,
with the possibility of alighting at will, and without being bound, on the single
point for a while. So in the ``aesthetic'' view of life depicted by
\textsc{S. Kierkegaard} in Part I of \textit{Enten-Eller}. Under all these forms
there plainly lies a regard for the whole of life in the background.

If a contradiction is experienced between what is momentary or sporadic and the
need for wholeness, there will arise a more or less conscious and a more or less
energetic striving to develop the personality into a work of art, in which every
single moment and every single capacity gets its place and comes into its own. No
element in the personal life is then regarded as a mere means, but always at the
same time as an end. What is a means or a point of transition for the development
of the personality shall so far as possible have immediate value itself. It is
\emph{the unforgettable fundamental thought of Greek (Platonic-Aristotelian)
ethics} concerning the harmonious unfolding of the life of the soul as the great
task --- a fundamental thought which has its significance quite particularly
under modern culture's tendency to isolate or to mechanize the several elements
of life. \textsc{Rousseau} and \textsc{Schiller} have renewed this ancient
fundamental thought --- against modern labour-helotism, and against dalliance and
decadence.

But the problem of continuity comes up again, and in a very sharp form, at the
question of the possibility that the single personality can form a completed
world for itself, and whether that would be the most valuable thing. Not only
must the individual always enter into relations of interaction with other
personalities in order to obtain means for his own development. There is also a
need for self-surrender, which can appear in different forms and which can lead
to ascribing immediate value to other personalities. The individual is thereby
drawn into the great realm of personalities; and just as the question before was
how far the several elements of life could be harmoniously ordered into the whole
of life of the single personality, so it now becomes a question how far the
several personalities can develop independently and yet in harmony with one
another, so that there can subsist \emph{a social whole of life} in analogy with
the individual whole of life. It is continuity within continuity that is here
striven towards. The measure of a society's perfection will --- in virtue of the
principle of continuity, which here shows itself plainly in its connection with
the principle of welfare --- be this: in what degree every personal being is
placed and treated in such a way that it does not stand as a mere means, but
always at the same time as an end. It is \textsc{Kant}'s famous proposition, with
a different grounding from the one he gave.\footnote{Compare my \textit{Etik}. 2,
p.~144--151.} Stoic and Christian ethics and the modern social and political
development lead in the same direction. The declaration of the ``rights of man''
(whether we regard it as a symptom or as a programme) rests on this
presupposition, and the sting in the social question has come about through this
presupposition not having been satisfied. In the discussion of more special
ethical and legal-philosophical questions too --- for example on the grounding of
monogamy and on the development of the penal system --- this principle furnishes
the measure. ---

Ethical work thus shows itself as a peculiar continuation of the great process of
existence. But the whole train of thought by which I have sought to ground this
point of view runs into difficulties which sharpen the ethical problem. We meet
here again the irrational.

\subsection*{b.\quad The Rationality of Ethical Valuation}
\label{sec:etiske-rationalitet}

\noindent\textbf{3.}~Ethics would be a more perfect science than it is if the
carrying through of the principle of continuity did not run into great
difficulties.

Every ethical argument holds only on the acknowledgement of a determinate
fundamental value which determines all special valuations. One may take one's
stand on the standpoint of the single moment or the single impulse, or on that of
the isolated personality, or on that of the family, or the class, or the state,
or humanity. The question is whether all such standpoints can really be brought
into harmony with one another, as has been assumed above; and the question is
especially whether it is possible, by way of argument, to bring a man who
consistently and with the requisite ruthlessness holds fast to one of these
standpoints over to one of the others. Between the fundamental value and the
special values a rational relation can perhaps be exhibited, so that one who
avows a fundamental value, and is sufficiently acquainted with the actual
conditions under which it is to be held fast, can be logically compelled to
concede the inferences one draws from the fundamental value. Thus self-assertion
and self-surrender have each their logic, and so too have the feeling for family
and the feeling for nation. But what when it is a question of \emph{transition
from one fundamental value to another}? Here inner consistency is not enough.
Consistency can only unfold and bring to consciousness what --- under certain
given, actual conditions --- follows from the fundamental value.
\textsc{Socrates} became the founder of ethics by his demand for self-knowledge,
which was in reality a demand to become clear about one's fundamental value and
about what follows consistently from it. But how the fundamental value is
procured, and whether there are not in fact different values which claim to be
fundamental values, he did not discuss more closely.

A fundamental value forms itself by a psychological and historical process which
presupposes more operative elements than logical consistency and knowledge of
facts. Through displacement of motive and of value\footnote{Compare my
\textit{Psykologi.} VI. B., 1--5. --- \textit{Etik} 2 XIII., 4.} one fundamental
value can pass over into another; but \emph{while} such a fundamental change is
going on in the domain of the life of feeling and will, it will be of no use to
argue from the presuppositions which only come to be present when the process is
finished. In education one cannot without more ado argue with the child from the
presuppositions of adults; the child has precisely first to acquire those
presuppositions. And so it goes too with the great process of education that goes
on in history. \emph{During} the education and the development there are of
course other fundamental motives at work than those which only become the result
of the whole process. He who is being educated will therefore not be able
rationally to understand the system by which he is being educated. ``Wer kann der
Raupe, die im Staube kriecht, von ihrem künftigen Futter sprechen?''

There does not, now, take place an evenly advancing education in history from
primitive fundamental values to higher ones (that is to say: to such as those who
have acquired them call higher in comparison with the primitive). History is the
great battlefield of fundamental values. Here individual stands against
individual, individual and society against each other, and society stands against
society, and it is often only after violent conflicts that a new fundamental
value can be fixed in the minds of men. An example of this is the way in which
the conquests of Alexander and of the Romans prepared the way for the development
of the general feeling of humanity. During this great struggle ethical thinking
can play a part only indirectly: by drawing the determinate consequences of each
single position, and by illuminating the significance of those consequences by as
many experiences as possible --- that is, by furthering self-knowledge in the
spirit of Socrates. In the struggle itself, or in the education, what is called
for is an art; science does not suffice, however significant a contribution it
may often make.

But we do not abandon the principle of continuity because we stand at the
boundary of scientific ethics. There, where it is not possible to pursue
continuity further in purely ethical form, we seek to trace it out in
psychological and historical form. Ethics passes at this point over into
psychology and into sociology.

While I have sought in my \textit{Etik} to show that justice, conceived as an
inner harmonious relation between self-assertion and self-surrender, and also
between feeling and thinking, is the highest quality of character, this has been
disputed by philosophical thinkers who maintain that in self-assertion and in
self-surrender such different tendencies (or, as I should say, fundamental
values) make themselves known that it is not possible to unite them in one
embracing concept.\footnote{\textsc{Francis Bradley}: \textit{Appearance and
Reality.} p.~414--418. --- \textsc{A. E. Taylor}: \textit{The Problem of
Conduct.} London 1901. p.~201.} And in popular literature self-surrender
(altruism) has been asserted just as ruthlessly by \textsc{Leo Tolstoy} as
self-assertion has by \textsc{Friedrich Nietzsche}. Here we have an example of a
strife between fundamental values that press forward at the same time. There will
probably always, in actual human life, be oscillations between these two opposite
poles. In the race as a whole both self-assertion and self-surrender have their
function to exercise. But the valuation of the single oscillation in the one or
the other direction will always depend on whether it leads forward in the
direction of an ordering of life within which every single personality can
develop as peculiarly and independently as possible, in such a way that precisely
thereby it renders the best help to the equally peculiar and independent
development of others. In individuals self-assertion and self-surrender will
stand in highly different relations to one another, and the harmony between the
two, in which on my view justice consists, will therefore appear with a highly
different timbre in different persons and at different times. But that too
belongs to the richness of life --- and to its continuity. The manifoldness of
the nuances conditions the connection and makes it different from a dead
uniformity, which has no value either from an ethical point of view or from any
other.

\noindent\textbf{4.}~Another difficulty, one connected with discontinuity, arises
from the fact that the inferences drawn under certain conditions from a
fundamental value have to be applied to \emph{individuals who come with highly
different inner and outer conditions} for satisfying the demand which such an
inference grounds. Aptitudes and impulses are not the same in all individuals,
either in kind or in degree. One and the same demand addressed to different
individuals therefore lays on them a highly different ethical work. The
starting-point and the initial speed are different. Some individuals are perhaps
already on the way to an involuntary performance of the content of the demand
before the demand becomes conscious to them; others must struggle laboriously to
lay the first foundation. The law, the demand, must therefore be set differently
for different persons if it is really to become the same for all. Each must be
taxed according to ability. A thoroughgoing individualization of the ethical
demand must be undertaken, if ethics itself is not to sin against the proposition
that personality is always an end and never merely a means. The ethical demand
shall not be an abstract or external command, but shall answer to the ethical
possibilities in the single personality and be able to develop them further.
Legislation and pedagogy cannot here be kept absolutely apart from one another.
But this makes ethical decisions difficult in the several cases.\footnote{Compare
my \textit{Etik}. 2, p.~70--73.} Here again the world --- the world of
personalities --- is great, and our brain is small. Ethical thinking cannot
formulate a law that can be applied without more ado in the manifoldness of life.
And yet we must assume that in each single case only one single decision can be
the right one (Eng ist die Welt, und das Gehirn ist weit). Our ethical thoughts
circle about the narrow way that leads to ethical truth, and work their way
forward to it through a constant struggle with the irrational, which they seek to
overcome by finding as great an approximation as possible. Here, as within the
problem of knowledge, we do not end with an absolute conclusion, but hear an
Excelsior! --- even where our thought has strained itself to the utmost.

\section*{B.\quad The Religious Problem}
\label{sec:religiose}

\noindent\textbf{5.}~The very fact that a religious problem exists already shows
the significance of continuity. For that religion has become problematic hangs
together with a division of labour, a differentiation, having set in within the
domain of the spiritual life. In religion's classical ages it stands as the one
concentrated form in which all the requirements of the life of the spirit are
satisfied; religion is as such not merely help and comfort, but is poetry,
morality and science, or it is at any rate able to take these domains up into
itself in an organic way. After these various domains have emancipated
themselves, and each develops according to its own laws, the question arises how
far the life of the spirit --- as regards value --- has preserved its continuity
through that transition from concentration to differentiation. That a
psychological and a historical continuity is present there is no reason to doubt;
but has not value been lost for human life by that transition?

The answers to this question sound highly different. Some are content that the
same dogmas are still taught as in the old days. From another side it is
maintained, on the contrary, that this dogmatic continuity means nothing: the
essential thing is whether life is led in the same way --- whether there is
ethical continuity --- whether ``the Christianity of the New Testament'' still
exists. And finally there are those who hold that with the division of labour in
cultural history the age of religion is past, and that this means no loss of
value but on the contrary a gain for the life of the spirit. Besides these
principal standpoints, which can themselves appear in different nuances, there is
a series of different and less characteristic standpoints.

The answering of the question depends of course on what one takes to be the
essential thing in religion. Continuity cannot appear in such features as
separate the different religions and religious standpoints from one another; one
must seek out the most deeply moving tendencies, which can reveal themselves
under highly different forms, and which may perhaps also continue to work after
the division of labour in cultural history has come into force. This is the task
of the philosophy of religion, and it seeks to solve it especially along two
roads. First, it institutes a comparison between religion and other sides of the
life of the spirit in order to find in which region of the life of the spirit
religion belongs; it seeks to determine religion's psychological place. And
secondly, it institutes a comparison between the most important historical forms
of religion in order to see whether the psychological determination found wins
confirmation in experience.\footnote{In my \textit{Religionsfilosofi} I have ---
for a more comprehensive illumination of the problem --- instituted, besides the
psychological-historical investigation, an epistemological and an ethical one as
well. Here, in this shorter presentation, I keep especially to the
psychological-historical investigation, and only indicate quite briefly, here and
there, the epistemological and the ethical points of view.}

In both investigations we shall have occasion to see that it is not only the
principle of continuity that leads to the posing of the religious problem, but
that the very concept of religion hangs very closely together with this
principle.

\noindent\textbf{6.}~The religious feeling can be placed in a clear and simple
relation to the other feelings when it is conceived as determined by the
experiences man makes concerning the fate of values in the world of reality. It
is thereby seen at once in its difference from and in its connection with the
other feelings. To every feeling there corresponds a value. The feeling of life,
the intellectual, the aesthetic and the ethical feeling thus make known different
kinds of value. The maintenance and development of life, truth, beauty and
goodness are values man can have a part in without religious feeling. But when it
shows itself that life, truth, beauty and goodness must struggle in order to
subsist in the world, there arises a peculiar feeling which is no longer
determined by those values in and for themselves, but by their subsistence, their
advance or their decline in existence, so far as man is able to survey it. The
experiences man makes about this will be able to give rise to a need to believe
that the values subsist, even where they no longer make themselves known in the
reality man can survey, or even if they no longer appear in the same forms as
hitherto. Man will be most inclined to assume that the value which subsists is
the same one he has hitherto had a part in; but he may develop to the point of
holding fast that the empirical values must perish in order that a higher and
more comprehensive connection of values may be satisfied. From men's
representations of a world of gods and of a life beyond, one comes to know which
values have lain nearest their hearts, and in what degree they have reconciled
themselves to the thought that the empirical values must undergo a more or less
essential metamorphosis if a subsistence of value is to be possible. In the end
what subsists is something that no eye has seen and no ear has heard. Religious
faith asserts a continuity in the domain of values beyond every possible
experience. This continuity of value may be of different content and extent. But
there will be a tendency to extend it beyond the human world and to conceive
existence as a home for the development and the subsistence of values. This
religious faith in continuity forms an analogy to the intellectual presupposition
of the rationality of existence (III), and the difficulties it has to struggle
with are analogous too.

Religious faith will have in particular a certain likeness to (and a certain
sympathy with) metaphysical idealism (III, 4). But it does not stand or fall with
it. The essential thing in religious faith is not the intellectual satisfaction
it can afford. The circle of representations in which religion is expressed does
not belong to its innermost psychological essence. What is central in religion is
--- if the psychological determination given is correct --- an interest of
feeling and of will. Intellectually we ask only about recognition (identity),
consequence (rationality) and causal connection (causality) (II, 1). Faith is
always only an \emph{object} for science, never itself science. It arises through
the harmonious or disharmonious relation between the picture of reality which
knowledge gives man, and the values which stand for man as the highest.

Possibly the difference we are compelled to make between value and reality holds
only from a human standpoint. But --- another standpoint we cannot, once for all,
take up.

The faith in the subsistence of value can itself acquire value by keeping courage
up in the struggle for life, and by spurring us to find new values as equivalents
for the values that disappear. Even one who assumes that the age of religion is
past --- an assumption which, if it is to be more than an assertion or a wish,
must be grounded epistemologically, psychologically and ethically --- even he
will surely feel a need to find equivalents for the loss of value suffered
through religion's disappearance. In this sense a religious problem exists also
for one who holds that religion's kernel disappears along with its shell. For
religion, by asserting the subsistence of value, acquires value itself, value in
the second power. And it designates in particular one of the most concentrated
forms of spiritual life that experience offers: it stands, where it is genuine
and original, as the collected result of man's life of feeling and will, and it
gets itself expression in the mightiest and noblest images that man's life of
representation commands. Probably this concentrated life-process will always
express itself under new forms --- whether or not one goes on using the word
religion.

\noindent\textbf{7.}~A confirmation of the characterization given of what is
essential in religion may be sought in the fact that the differences between the
religions, or between the religious standpoints among themselves, can simply and
easily be explained from it. These differences rest partly on differences in the
values in whose subsistence there is belief, partly on differences in the
conception of reality that lies at the base, partly on differences in experience
concerning the relation between the values and reality.

The differences with respect to \emph{the values} appear in the historical
religions in what is regarded as the task of the gods --- what one thinks of the
gods as struggling for --- or, to use an expression of \textsc{Plato}'s, in
\emph{that which makes God divine}.\footnote{\textsc{Plato}: \textit{Phaedrus.}
p.~249 C (πρὸς οἷσπερ θεὸς ὢν θεῖός ἐστι). (More exactly translated: ``that in
which it is God's divinity to live''.)} It may be the purely physical maintenance
of life and the enjoyments attached to it, or it may be the good, the beautiful
and the true. Historians of religion emphasize the transition from
nature-religion to ethical religion as the most significant within the history of
religion; but this transition rests precisely on a transition from the more
elementary values to the ideal values as that which lies at the base of the
relation to the gods --- what the gods guard, and to which they themselves owe
their divinity. At the same time the more external relation between the powers
man believes himself to stand in relation to, and that which those powers guard,
falls away, and the gods themselves become more immediate representatives of the
valuable which they uphold --- indeed become one with it. An ever more inward
connection between the religious and the ethical can be traced in religion's
development, and religion tends more and more to become a projection of the
ethical.

When \emph{the picture of reality} is altered by advancing knowledge, there
occurs a no less decisive alteration in religion's character than when the values
lying at the base are altered. The most conspicuous religious crises are due to a
new conception of reality. Religion has here on the whole behaved receptively
towards the knowledge developed along another road, but has afterwards sought to
take it up and to use it as a foundation for its images. Thus animism, the
knowledge of the stars, the doctrine of the transmigration of souls, Indian and
Greek philosophy, Copernicanism and more recent natural science, speculative and
critical philosophy have successively intervened --- now as welcome forms for the
religious life of representation, now as widening impulses, now as hostile
tendencies.

But the decisive thing will always be \emph{the relation between the values and
reality}. According as this relation appears in man's experience of life as
harmonious or as disharmonious, religion will acquire a different character; and
in the highest religions both life's disharmonies and its harmonies will find
their place --- there the religious mood will first reach height and firmness
when it has worked its way through conflict and suffering to blessedness. The
strength of the disharmonies then becomes only a testimony to the strength of the
victorious harmony. But to this comes another relation. Religion will acquire a
different character according as what is highest in value is thought as
subsisting eternally, raised above all becoming and change, so that life in time
is in the end an illusion, or as the valuable itself unfolding in time and
struggling for its subsistence in the course of the changes. On this difference
rests the opposition between the Indian-Greek type of religion and the
Persian-Christian type.

It seems then that it is no alien measure we apply to the religions when we
characterize and value them according to the way in which, and the degree to
which, the principle of the subsistence of value appears in them. It is a measure
that naturally presents itself for a comparative consideration of the history of
the religions and of religious standpoints.

The second of the types of religion named stands in history as the victorious
one. It has certainly not reached its definitive form. If it comes to take on new
forms, that will be due, as in the past, to prophetic personalities. They alone
are able to weave the Godhead a new garment.

For the philosopher the essential interest is to exhibit that so long as a
religious problem subsists, it hangs closely together partly with the need for
continuity, partly with the time-relation, partly with the ethical values. The
kinship and the inner connection between the great problems of the human spirit,
the practical as well as the theoretical, here come plainly forward.

And at all the problems we end --- when we see the opposition in its sharpest
form --- in the eternal strife, a strife in which the spiritual powers of men
shall ever anew be tried. But because we cannot solve the great problems, we can
nevertheless see the way that leads forward in such a manner that the right both
of thought and of life is maintained. The insolubility of the problems means
really only this: that however far we get in our inquiry and our thinking, there
form themselves for us ever new horizons, new goals and new tasks.

\end{document}
